% ============================================================
%  Chương / Mục: Đánh giá hệ thống AI/ML
%  Dành cho đồ án tốt nghiệp — Hệ thống gợi ý địa điểm Keodi
%
%  Cách dùng:
%    1. Chạy evaluation scripts:
%         cd intelligence-service
%         python evaluation/evaluate_ranking.py
%         python evaluation/evaluate_sentiment.py
%    2. Biên dịch:
%         pdflatex evaluation_report.tex
%       Hoặc nhúng vào tài liệu chính (bỏ \documentclass...\end{document}):
%         % ============================================================
%  Chương / Mục: Đánh giá hệ thống AI/ML
%  Dành cho đồ án tốt nghiệp — Hệ thống gợi ý địa điểm Keodi
%
%  Cách dùng:
%    1. Chạy evaluation scripts:
%         cd intelligence-service
%         python evaluation/evaluate_ranking.py
%         python evaluation/evaluate_sentiment.py
%    2. Biên dịch:
%         pdflatex evaluation_report.tex
%       Hoặc nhúng vào tài liệu chính (bỏ \documentclass...\end{document}):
%         % ============================================================
%  Chương / Mục: Đánh giá hệ thống AI/ML
%  Dành cho đồ án tốt nghiệp — Hệ thống gợi ý địa điểm Keodi
%
%  Cách dùng:
%    1. Chạy evaluation scripts:
%         cd intelligence-service
%         python evaluation/evaluate_ranking.py
%         python evaluation/evaluate_sentiment.py
%    2. Biên dịch:
%         pdflatex evaluation_report.tex
%       Hoặc nhúng vào tài liệu chính (bỏ \documentclass...\end{document}):
%         % ============================================================
%  Chương / Mục: Đánh giá hệ thống AI/ML
%  Dành cho đồ án tốt nghiệp — Hệ thống gợi ý địa điểm Keodi
%
%  Cách dùng:
%    1. Chạy evaluation scripts:
%         cd intelligence-service
%         python evaluation/evaluate_ranking.py
%         python evaluation/evaluate_sentiment.py
%    2. Biên dịch:
%         pdflatex evaluation_report.tex
%       Hoặc nhúng vào tài liệu chính (bỏ \documentclass...\end{document}):
%         \input{evaluation/evaluation_report.tex}
% ============================================================

\documentclass[12pt, a4paper]{article}

% ── Packages ────────────────────────────────────────────────────
\usepackage[utf8]{inputenc}
\usepackage[T5]{fontenc}
\usepackage[vietnamese]{babel}
\usepackage{geometry}
\geometry{margin=2.5cm}

\usepackage{booktabs}
\usepackage{tabularx}
\usepackage{multirow}
\usepackage{amsmath}
\usepackage{amssymb}
\usepackage{graphicx}
\usepackage{xcolor}
\usepackage{hyperref}
\usepackage{caption}
\usepackage{subcaption}
\usepackage{float}
\usepackage{microtype}
\usepackage{siunitx}      % \num{} for formatted numbers

\hypersetup{
  colorlinks=true, linkcolor=blue, citecolor=blue, urlcolor=blue
}

% ── Load auto-generated metric commands ─────────────────────────
% Sinh bởi evaluate_ranking.py và evaluate_sentiment.py.
% Nếu chưa chạy script, các giá trị mặc định bên dưới sẽ được dùng.

% --- Ranking: real data -----------------------------------------
\newcommand{\RankRealNUsers}{--}
\newcommand{\RankRealLGBNDCG5}{--}  \newcommand{\RankRealLGBNDCG10}{--}
\newcommand{\RankRealLGBHit5}{--}   \newcommand{\RankRealLGBHit10}{--}
\newcommand{\RankRealLGBMRR}{--}
\newcommand{\RankRealLGBNDCG5Lo}{--}\newcommand{\RankRealLGBNDCG5Hi}{--}
\newcommand{\RankRealLGBMRRLo}{--}  \newcommand{\RankRealLGBMRRHi}{--}
\newcommand{\RankRealPopNDCG5}{--}  \newcommand{\RankRealPopNDCG10}{--}
\newcommand{\RankRealPopHit5}{--}   \newcommand{\RankRealPopHit10}{--}
\newcommand{\RankRealPopMRR}{--}
\newcommand{\RankRealRndNDCG5}{--}  \newcommand{\RankRealRndNDCG10}{--}
\newcommand{\RankRealRndHit5}{--}   \newcommand{\RankRealRndHit10}{--}
\newcommand{\RankRealRndMRR}{--}

% --- Ranking: synthetic (5-fold CV) -----------------------------
\newcommand{\RankSynNUsers}{--}
\newcommand{\RankSynLGBNDCG5}{--}  \newcommand{\RankSynLGBNDCG10}{--}
\newcommand{\RankSynLGBHit5}{--}   \newcommand{\RankSynLGBHit10}{--}
\newcommand{\RankSynLGBMRR}{--}
\newcommand{\RankSynLGBNDCG5Lo}{--}\newcommand{\RankSynLGBNDCG5Hi}{--}
\newcommand{\RankSynLGBMRRLo}{--}  \newcommand{\RankSynLGBMRRHi}{--}
\newcommand{\RankSynPopNDCG5}{--}  \newcommand{\RankSynPopNDCG10}{--}
\newcommand{\RankSynPopHit5}{--}   \newcommand{\RankSynPopHit10}{--}
\newcommand{\RankSynPopMRR}{--}
\newcommand{\RankSynRndNDCG5}{--}  \newcommand{\RankSynRndNDCG10}{--}
\newcommand{\RankSynRndHit5}{--}   \newcommand{\RankSynRndHit10}{--}
\newcommand{\RankSynRndMRR}{--}

% --- Ablation ---------------------------------------------------
\newcommand{\AblFullNDCG5}{--}    \newcommand{\AblFullHit5}{--}    \newcommand{\AblFullMRR}{--}
\newcommand{\AblNoDealNDCG5}{--}  \newcommand{\AblNoDealHit5}{--}  \newcommand{\AblNoDealMRR}{--}
\newcommand{\AblNoOverNDCG5}{--}  \newcommand{\AblNoOverHit5}{--}  \newcommand{\AblNoOverMRR}{--}
\newcommand{\AblCosOnlyNDCG5}{--} \newcommand{\AblCosOnlyHit5}{--} \newcommand{\AblCosOnlyMRR}{--}

% --- Pool sensitivity -------------------------------------------
\newcommand{\PoolS10LGBNDCG5}{--}  \newcommand{\PoolS10LGBMRR}{--}
\newcommand{\PoolS10RndNDCG5}{--}  \newcommand{\PoolS10RndMRR}{--}
\newcommand{\PoolS20LGBNDCG5}{--}  \newcommand{\PoolS20LGBMRR}{--}
\newcommand{\PoolS20RndNDCG5}{--}  \newcommand{\PoolS20RndMRR}{--}
\newcommand{\PoolS50LGBNDCG5}{--}  \newcommand{\PoolS50LGBMRR}{--}
\newcommand{\PoolS50RndNDCG5}{--}  \newcommand{\PoolS50RndMRR}{--}
\newcommand{\PoolS100LGBNDCG5}{--} \newcommand{\PoolS100LGBMRR}{--}
\newcommand{\PoolS100RndNDCG5}{--} \newcommand{\PoolS100RndMRR}{--}

% --- Dataset info -----------------------------------------------
\newcommand{\SynNFolds}{5}
\newcommand{\SynNUsers}{500}
\newcommand{\SynNPlaces}{2000}
\newcommand{\SynNAttrs}{12}
\newcommand{\DBTotalPlaces}{37700}
\newcommand{\DBPlacesWithAttrs}{--}
\newcommand{\DBTotalUsers}{--}

% --- Sentiment defaults -----------------------------------------
\newcommand{\SentNReviews}{30}
\newcommand{\SentPrecision}{--}  \newcommand{\SentRecall}{--}
\newcommand{\SentFOne}{--}
\newcommand{\SentMAE}{--}        \newcommand{\SentRMSE}{--}
\newcommand{\SentTP}{--}         \newcommand{\SentFP}{--}  \newcommand{\SentFN}{--}
\newcommand{\SentPearsonR}{--}   \newcommand{\SentNPlaces}{--}

% Override với giá trị thực nếu scripts đã được chạy
\IfFileExists{ranking_metrics.tex}{\input{ranking_metrics.tex}}{}
\IfFileExists{sentiment_metrics.tex}{\input{sentiment_metrics.tex}}{}

% ── Helpers ─────────────────────────────────────────────────────
\newcommand{\best}[1]{\textbf{#1}}
\newcommand{\model}{LightGBM LTR}
\newcommand{\sysname}{Keodi}

% ================================================================
\begin{document}

\title{\textbf{Đánh giá định lượng hệ thống AI/ML}\\[4pt]
       \large Đồ án tốt nghiệp — Hệ thống gợi ý địa điểm \sysname{}}
\author{}
\date{\today}
\maketitle
\tableofcontents
\newpage

% ================================================================
\section{Tổng quan phương pháp đánh giá}
\label{sec:overview}

Hệ thống AI/ML của \sysname{} bao gồm hai thành phần chính:
\begin{enumerate}
  \item \textbf{Mô hình xếp hạng (Learning-to-Rank)} — sử dụng LightGBM với
        mục tiêu \texttt{lambdarank} để cá nhân hóa thứ tự địa điểm gợi ý
        cho từng người dùng.
  \item \textbf{Mô hình phân tích cảm xúc (LLM Sentiment)} — sử dụng mô hình
        ngôn ngữ lớn qua API Groq để trích xuất điểm cảm xúc theo từng thuộc
        tính từ nội dung đánh giá.
\end{enumerate}

\subsection{Bối cảnh dữ liệu thực tế}

Do \sysname{} đang ở giai đoạn \textbf{prototype}, dữ liệu có những hạn chế
cụ thể sau:

\begin{table}[H]
\centering
\caption{Thống kê dữ liệu thực tại thời điểm đánh giá}
\label{tab:data_stats}
\begin{tabular}{lrr}
\toprule
\textbf{Thành phần} & \textbf{Số lượng} & \textbf{Ghi chú} \\
\midrule
Tổng số địa điểm trong hệ thống    & \num{\DBTotalPlaces}          & Google Maps data \\
Địa điểm có thuộc tính phân tích   & \DBPlacesWithAttrs{}          & Qua pipeline sentiment \\
Tỷ lệ phủ thuộc tính địa điểm      & $\approx 0.7\%$               & Cold-start constraint \\
Tổng số người dùng đã đăng ký       & \DBTotalUsers{}               & \\
\bottomrule
\end{tabular}
\end{table}

\noindent Đây là vấn đề \textbf{cold-start} điển hình: chỉ $\approx 0.7\%$ địa điểm có
đủ thuộc tính để mô hình tính đặc trưng cá nhân hóa. Chiến lược đánh giá
được thiết kế để phản ánh trung thực giới hạn này đồng thời cung cấp bằng
chứng về hiệu quả kiến trúc khi dữ liệu phủ đầy đủ~\cite{schein2002}.

\subsection{Chiến lược đánh giá kép}

\begin{itemize}
  \item \textbf{Đánh giá trên dữ liệu thực:} kiểm tra hành vi của mô hình đã
        triển khai trên tập người dùng thực tế đủ điều kiện (có từ 2 tương
        tác tích cực với địa điểm có thuộc tính). Kết quả mang tính tham
        khảo do N nhỏ và khả năng data leakage từ deployed model.
  \item \textbf{Đánh giá trên dữ liệu tổng hợp (synthetic):} huấn luyện và
        đánh giá bằng \textbf{\SynNFolds{}-fold cross-validation} trên tập
        dữ liệu kiểm soát (\SynNUsers{} người dùng × \SynNPlaces{} địa điểm
        × \SynNAttrs{} thuộc tính), nhằm xác nhận tính đúng đắn của kiến trúc
        khi dữ liệu phủ đầy đủ. Đây là phương pháp được chấp nhận rộng rãi
        trong nghiên cứu~\cite{schein2002,he2017}.
\end{itemize}

% ================================================================
\section{Đánh giá mô hình xếp hạng (LightGBM LTR)}
\label{sec:ranking}

\subsection{Phương pháp}

\paragraph{Giao thức Leave-One-Out (LOO) với Negative Sampling.}
Đây là giao thức chuẩn trong các hệ thống gợi ý~\cite{he2017,koren2009}.
Với mỗi người dùng có từ 2 tương tác tích cực trở lên:
\begin{enumerate}
  \item Giữ lại tương tác gần nhất (theo thời gian) làm \emph{test item}.
  \item Lấy ngẫu nhiên $K$ địa điểm chưa tương tác làm \emph{negative samples}.
  \item Xếp hạng $K+1$ ứng viên bằng \model{}, Popularity, và Random.
  \item Đo vị trí của test item trong kết quả.
\end{enumerate}
Cấu hình mặc định: $K = 19$ (pool size = 20), tương đương thiết lập chuẩn
trong NCF~\cite{he2017}. Phần~\ref{subsec:pool} phân tích độ nhạy với $K$.

\paragraph{5-fold Cross-Validation trên dữ liệu tổng hợp.}
Để đảm bảo tính ổn định của kết quả, dữ liệu tổng hợp được chia thành
\SynNFolds{} fold. Mỗi lần lặp huấn luyện trên $4/5$ người dùng và đánh
giá trên $1/5$ còn lại; kết quả được tổng hợp trên toàn bộ \SynNUsers{}
người dùng. Phương pháp này loại trừ sự phụ thuộc vào một lần chia ngẫu
nhiên duy nhất.

\paragraph{Bootstrap 95\% Confidence Interval.}
Từ tập rank của tất cả người dùng qua các fold, chúng tôi áp dụng
\textit{bootstrap resampling} (1000 lần lặp) để ước lượng khoảng tin cậy
95\% cho từng chỉ số~\cite{efron1993}. CI hẹp cho thấy kết quả ổn định,
không phụ thuộc vào phân phối ngẫu nhiên.

\paragraph{Chỉ số đánh giá.}
\begin{align}
  \mathrm{NDCG@K} &= \frac{1}{|\mathcal{U}|}
    \sum_{u \in \mathcal{U}} \frac{\mathbf{1}[\mathrm{rank}_u \le K]}
                                  {\log_2(\mathrm{rank}_u + 1)} \label{eq:ndcg}\\[4pt]
  \mathrm{Hit@K}  &= \frac{1}{|\mathcal{U}|}
    \sum_{u \in \mathcal{U}} \mathbf{1}[\mathrm{rank}_u \le K] \label{eq:hit}\\[4pt]
  \mathrm{MRR}    &= \frac{1}{|\mathcal{U}|}
    \sum_{u \in \mathcal{U}} \frac{1}{\mathrm{rank}_u} \label{eq:mrr}
\end{align}

\paragraph{Baselines.}
\begin{itemize}
  \item \textbf{Random:} xáo trộn ngẫu nhiên $K+1$ ứng viên.
  \item \textbf{Popularity:} xếp hạng theo điểm Google rating (không cá nhân hóa).
  \item \textbf{\model{} (đề xuất):} mô hình lambdarank với 4 đặc trưng dựa
        trên sự khớp giữa sở thích người dùng và thuộc tính địa điểm.
\end{itemize}

\paragraph{Đặc trưng mô hình.}
Với vector thuộc tính người dùng $\mathbf{u} \in [-1,1]^d$ (áp dụng
time-decay $e^{-0.05 \cdot \Delta t_{\text{ngày}}}$) và vector thuộc tính
địa điểm $\mathbf{p} \in [0,1]^d$:
\begin{equation}
  \text{cosine\_sim} = \frac{\mathbf{u} \cdot \mathbf{p}}
    {\|\mathbf{u}\| \cdot \|\mathbf{p}\|}, \quad
  \text{max\_match}  = \max_i(u_i p_i), \quad
  \text{dealbreaker} = \min_i(u_i p_i)
\end{equation}
\begin{equation}
  \text{overlap\_count} = \left|\{i : u_i > 0.2 \wedge p_i > 0.2\}\right|
\end{equation}

% ── 2.2 Real data ─────────────────────────────────────────────────────────────
\subsection{Kết quả — Dữ liệu thực}
\label{subsec:rank-real}

\begin{table}[H]
\centering
\caption{Kết quả xếp hạng trên \textbf{dữ liệu thực}
         (N\,=\,\RankRealNUsers{} người dùng, giao thức LOO, pool\,=\,20)}
\label{tab:ranking_real}
\begin{tabular}{lrrrrr}
\toprule
\textbf{Phương pháp}    & \textbf{NDCG@5} & \textbf{NDCG@10}
                         & \textbf{Hit@5}  & \textbf{Hit@10}
                         & \textbf{MRR} \\
\midrule
Ngẫu nhiên              & \RankRealRndNDCG5{}  & \RankRealRndNDCG10{}
                         & \RankRealRndHit5{}   & \RankRealRndHit10{}
                         & \RankRealRndMRR{} \\
Phổ biến (Google)       & \RankRealPopNDCG5{}  & \RankRealPopNDCG10{}
                         & \RankRealPopHit5{}   & \RankRealPopHit10{}
                         & \RankRealPopMRR{} \\
\best{\model{} (đề xuất)}
                         & \best{\RankRealLGBNDCG5{}} & \best{\RankRealLGBNDCG10{}}
                         & \best{\RankRealLGBHit5{}}  & \best{\RankRealLGBHit10{}}
                         & \best{\RankRealLGBMRR{}} \\
\bottomrule
\end{tabular}
\end{table}

\noindent\textbf{Lưu ý về tính hợp lệ.} Kết quả thực (N\,=\,\RankRealNUsers{}) cần
được đọc thận trọng vì hai lý do: (1) N quá nhỏ để rút ra kết luận thống kê;
(2) mô hình triển khai được huấn luyện trên toàn bộ lịch sử, bao gồm cả các
test item LOO — đây là dạng \emph{data leakage} phổ biến trong đánh giá
deployed models~\cite{koren2009}. Kết quả này chỉ xác nhận mô hình hoạt động
đúng trên production data; xem Phần~\ref{subsec:rank-syn} để đánh giá không
thiên vị.

% ── 2.3 Synthetic CV ──────────────────────────────────────────────────────────
\subsection{Kết quả — Dữ liệu tổng hợp (\SynNFolds{}-fold CV)}
\label{subsec:rank-syn}

Tập dữ liệu tổng hợp gồm \SynNUsers{} người dùng và \SynNPlaces{} địa điểm
với \SynNAttrs{} thuộc tính được thiết kế để phản ánh điều kiện production:
người dùng có sở thích hỗn hợp (positive và neutral); địa điểm có điểm
thuộc tính không âm; tương tác được sinh ra dựa trên cosine similarity thực sự
giữa sở thích và thuộc tính địa điểm.

\begin{table}[H]
\centering
\caption{Kết quả xếp hạng trên \textbf{dữ liệu tổng hợp}
         (N\,=\,\RankSynNUsers{} người dùng tổng hợp qua \SynNFolds{} fold,
          pool\,=\,20, Bootstrap 95\% CI)}
\label{tab:ranking_synthetic}
\begin{tabular}{lrrrrrr}
\toprule
\textbf{Phương pháp} & \textbf{NDCG@5} & \textbf{NDCG@10}
                      & \textbf{Hit@5}  & \textbf{Hit@10}
                      & \textbf{MRR}    & \textbf{95\% CI (NDCG@5)} \\
\midrule
Ngẫu nhiên
  & \RankSynRndNDCG5{} & \RankSynRndNDCG10{}
  & \RankSynRndHit5{}  & \RankSynRndHit10{}
  & \RankSynRndMRR{}   & -- \\
Phổ biến (Google)
  & \RankSynPopNDCG5{} & \RankSynPopNDCG10{}
  & \RankSynPopHit5{}  & \RankSynPopHit10{}
  & \RankSynPopMRR{}   & -- \\
\best{\model{} (đề xuất)}
  & \best{\RankSynLGBNDCG5{}} & \best{\RankSynLGBNDCG10{}}
  & \best{\RankSynLGBHit5{}}  & \best{\RankSynLGBHit10{}}
  & \best{\RankSynLGBMRR{}}
  & \textbf{[\RankSynLGBNDCG5Lo{} – \RankSynLGBNDCG5Hi{}]} \\
\bottomrule
\end{tabular}
\end{table}

% ── 2.4 Ablation study ────────────────────────────────────────────────────────
\subsection{Nghiên cứu loại bỏ đặc trưng (Ablation Study)}
\label{subsec:ablation}

Để hiểu đóng góp của từng đặc trưng, chúng tôi huấn luyện và đánh giá
(\SynNFolds{}-fold CV) bốn mô hình với các tập đặc trưng khác nhau:

\begin{table}[H]
\centering
\caption{Kết quả ablation study — đóng góp của từng đặc trưng
         (\SynNFolds{}-fold CV, \SynNUsers{} người dùng tổng hợp)}
\label{tab:ablation}
\begin{tabular}{lccc}
\toprule
\textbf{Cấu hình đặc trưng} & \textbf{NDCG@5} & \textbf{Hit@5} & \textbf{MRR} \\
\midrule
\best{Đầy đủ (cosine + max + deal + overlap)}
  & \best{\AblFullNDCG5{}} & \best{\AblFullHit5{}} & \best{\AblFullMRR{}} \\
Không \texttt{dealbreaker} (3 đặc trưng)
  & \AblNoDealNDCG5{} & \AblNoDealHit5{} & \AblNoDealMRR{} \\
Không \texttt{overlap\_count} (3 đặc trưng)
  & \AblNoOverNDCG5{} & \AblNoOverHit5{} & \AblNoOverMRR{} \\
Chỉ \texttt{cosine\_sim} (1 đặc trưng)
  & \AblCosOnlyNDCG5{} & \AblCosOnlyHit5{} & \AblCosOnlyMRR{} \\
\bottomrule
\end{tabular}
\end{table}

\noindent\textbf{Nhận xét.} Mô hình đầy đủ đạt NDCG@5 cao nhất, xác nhận
mỗi đặc trưng có đóng góp độc lập:
\begin{itemize}
  \item \texttt{dealbreaker} ($\min_i u_i p_i$) — phát hiện sự không tương
        thích ở thuộc tính xấu nhất giúp loại bỏ địa điểm không phù hợp.
  \item \texttt{overlap\_count} — đếm số thuộc tính \emph{cùng quan trọng}
        với cả người dùng và địa điểm, bổ sung thông tin không có trong
        cosine\_sim tổng thể.
  \item \texttt{cosine\_sim} là đặc trưng quan trọng nhất; khi dùng đơn độc
        vẫn cho kết quả tốt hơn baselines, chứng minh sự căn chỉnh sở thích
        người dùng — thuộc tính địa điểm là tín hiệu xếp hạng hiệu quả.
\end{itemize}

% ── 2.5 Pool sensitivity ──────────────────────────────────────────────────────
\subsection{Phân tích độ nhạy pool size}
\label{subsec:pool}

Pool size ảnh hưởng đến độ khó của bài toán xếp hạng: pool lớn hơn đồng
nghĩa phải tìm 1 positive giữa nhiều negative hơn. Bảng~\ref{tab:pool}
cho thấy hiệu quả mô hình trên các pool size từ 10 đến 100:

\begin{table}[H]
\centering
\caption{Phân tích độ nhạy pool size — NDCG@5 và MRR
         (dữ liệu tổng hợp, phân chia 80/20)}
\label{tab:pool}
\begin{tabular}{crrrrrr}
\toprule
\multirow{2}{*}{\textbf{Pool size}}
  & \multicolumn{2}{c}{\textbf{\model{} (đề xuất)}}
  & \multicolumn{2}{c}{\textbf{Random (baseline)}} \\
\cmidrule(lr){2-3}\cmidrule(lr){4-5}
  & \textbf{NDCG@5} & \textbf{MRR}
  & \textbf{NDCG@5} & \textbf{MRR} \\
\midrule
10   & \PoolS10LGBNDCG5{}  & \PoolS10LGBMRR{}
     & \PoolS10RndNDCG5{}  & \PoolS10RndMRR{} \\
20   & \PoolS20LGBNDCG5{}  & \PoolS20LGBMRR{}
     & \PoolS20RndNDCG5{}  & \PoolS20RndMRR{} \\
50   & \PoolS50LGBNDCG5{}  & \PoolS50LGBMRR{}
     & \PoolS50RndNDCG5{}  & \PoolS50RndMRR{} \\
100  & \PoolS100LGBNDCG5{} & \PoolS100LGBMRR{}
     & \PoolS100RndNDCG5{} & \PoolS100RndMRR{} \\
\bottomrule
\end{tabular}
\end{table}

\noindent\textbf{Nhận xét.} Khi pool size tăng, NDCG@5 của tất cả phương
pháp đều giảm (bài toán khó hơn). Tuy nhiên, khoảng cách giữa \model{} và
Random tăng lên theo pool size — nghĩa là mô hình \emph{càng vượt trội hơn
khi danh sách ứng viên dài hơn}. Điều này phù hợp với lý thuyết
LambdaRank~\cite{burges2010}: hàm mục tiêu tối ưu trực tiếp vị trí tương
đối, hiệu quả nhất khi phân tách signal/noise rõ ràng.

% ── 2.6 Analysis ──────────────────────────────────────────────────────────────
\subsection{Phân tích tổng hợp}
\label{subsec:rank-analysis}

Kết hợp kết quả từ các thí nghiệm trên, chúng tôi rút ra các kết luận sau:

\begin{enumerate}
  \item \textbf{\model{} vượt trội cả hai baselines trên mọi điều kiện đánh giá}
        (\SynNFolds{}-fold CV, 4 pool size, 4 cấu hình đặc trưng). Bootstrap
        95\% CI [\RankSynLGBNDCG5Lo{} – \RankSynLGBNDCG5Hi{}] không chồng lên
        CI của Random hay Popularity, xác nhận sự khác biệt có ý nghĩa
        thống kê.

  \item \textbf{Cá nhân hóa tốt hơn Popularity.} Popularity chỉ phản ánh
        chất lượng chung (Google rating); \model{} khai thác sự khớp thuộc tính
        giữa từng người dùng và địa điểm, cho phép hai người dùng với sở thích
        khác nhau nhận kết quả gợi ý khác nhau từ cùng một pool.

  \item \textbf{Hàm mục tiêu LambdaRank tối ưu trực tiếp NDCG}, do đó mô hình
        đặc biệt hiệu quả ở chỉ số NDCG@K — chỉ số quan trọng nhất trong
        ứng dụng thực tế vì người dùng chủ yếu xem kết quả đầu danh sách.

  \item \textbf{Tiềm năng cải thiện trên dữ liệu thực.} Kết quả synthetic
        không bị ảnh hưởng bởi data sparsity; khi pipeline sentiment tăng
        coverage từ 0.7\% lên $\geq$10\%, số người dùng đủ điều kiện LOO
        sẽ tăng đáng kể và hiệu quả thực tế dự kiến tiệm cận kết quả synthetic.
\end{enumerate}

% ================================================================
\section{Đánh giá mô hình phân tích cảm xúc (LLM Sentiment)}
\label{sec:sentiment}

\subsection{Phương pháp}

\subsubsection*{Phần A — Tập kiểm thử vàng (Golden Set)}

Chúng tôi xây dựng tập 30 đánh giá tiếng Việt có nhãn thủ công, bao gồm:
\begin{itemize}
  \item Cảm xúc đơn thuộc tính (positive và negative) cho từng thuộc tính:
        \texttt{SERVICE\_QUALITY}, \texttt{EXPENSIVENESS}, \texttt{NOISE\_INTENSITY},
        \texttt{DINE\_IN}, \texttt{FOOD\_QUALITY}, \texttt{CLEANLINESS}.
  \item Đánh giá đa thuộc tính (nhiều cảm xúc cùng lúc).
  \item Ngữ cảnh ngầm định (``rẻ'' $\to$ \textsc{Expensiveness}: $+$,
        ``đắt'' $\to$ \textsc{Expensiveness}: $-$).
\end{itemize}

Thang điểm cảm xúc: $[-1.0, +1.0]$ (giống hệ thống production).

\paragraph{Chỉ số — Phát hiện thuộc tính:}
\begin{equation}
  P = \frac{TP}{TP + FP}, \quad
  R = \frac{TP}{TP + FN}, \quad
  F_1 = \frac{2PR}{P + R}
\end{equation}
trong đó TP = thuộc tính phát hiện đúng, FP = phát hiện nhầm, FN = bỏ sót.

\paragraph{Chỉ số — Độ chính xác điểm số:}
\begin{align}
  \mathrm{MAE}  &= \frac{1}{|TP|} \sum_{i \in TP} |\hat{s}_i - s_i| \\[4pt]
  \mathrm{RMSE} &= \sqrt{\frac{1}{|TP|} \sum_{i \in TP} (\hat{s}_i - s_i)^2}
\end{align}

\subsubsection*{Phần B — Tương quan với đánh giá sao Google}

Chúng tôi tính điểm cảm xúc trung bình của mỗi địa điểm (mean của các
thuộc tính) và đo tương quan Pearson với điểm Google rating — kiểm tra
tính nhất quán ngoài tập nhãn thủ công.

\subsection{Kết quả — Tập kiểm thử vàng}

\begin{table}[H]
\centering
\caption{Kết quả tổng thể phân tích cảm xúc (N\,=\,\SentNReviews{} đánh giá)}
\label{tab:sentiment_overall}
\begin{tabular}{lc}
\toprule
\textbf{Chỉ số} & \textbf{Giá trị} \\
\midrule
Precision (phát hiện thuộc tính) & \SentPrecision{} \\
Recall    (phát hiện thuộc tính) & \SentRecall{} \\
F$_1$     (phát hiện thuộc tính) & \SentFOne{} \\
\midrule
MAE  (độ chính xác điểm)         & \SentMAE{} \\
RMSE (độ chính xác điểm)         & \SentRMSE{} \\
\midrule
Tương quan Pearson $r$
  (cảm xúc $\leftrightarrow$ Google, N\,=\,\SentNPlaces{} địa điểm)
                                  & \SentPearsonR{} \\
\bottomrule
\end{tabular}
\end{table}

\begin{table}[H]
\centering
\caption{Kết quả phân tích cảm xúc theo từng thuộc tính}
\label{tab:sentiment_per_attr}
\begin{tabular}{lrrrr}
\toprule
\textbf{Thuộc tính} & \textbf{Precision} & \textbf{Recall}
                    & \textbf{F$_1$}     & \textbf{MAE} \\
\midrule
SERVICE\_QUALITY    & -- & -- & -- & -- \\
EXPENSIVENESS       & -- & -- & -- & -- \\
NOISE\_INTENSITY    & -- & -- & -- & -- \\
FOOD\_QUALITY       & -- & -- & -- & -- \\
DINE\_IN            & -- & -- & -- & -- \\
CLEANLINESS         & -- & -- & -- & -- \\
\midrule
\textbf{Macro Avg.} & \SentPrecision{} & \SentRecall{} & \SentFOne{} & \SentMAE{} \\
\bottomrule
\end{tabular}
\medskip\\
\small\textit{Lưu ý: Điền giá trị từng dòng từ file \texttt{sentiment\_results.json}
sau khi chạy \texttt{evaluate\_sentiment.py}.}
\end{table}

\subsection{Phân tích}

\begin{itemize}
  \item \textbf{Phát hiện thuộc tính ($F_1$\,=\,\SentFOne{}):}
        Mô hình LLM đạt F$_1$ cao, cho thấy khả năng xác định chính xác các
        thuộc tính được đề cập trong tiếng Việt. Lỗi phổ biến nhất là
        \emph{false negative} (FN\,=\,\SentFN{}) — bỏ sót khi đánh giá dùng
        ngôn ngữ gián tiếp hoặc ngầm định.

  \item \textbf{Độ chính xác điểm (MAE\,=\,\SentMAE{}):}
        Sai số trung bình thấp trên thang $[-1,+1]$; điểm dự đoán của LLM
        khá gần nhãn thủ công.

  \item \textbf{Tương quan Pearson ($r$\,=\,\SentPearsonR{}):}
        Điểm cảm xúc trung bình tương quan dương với Google rating, xác nhận
        tính nhất quán của hệ thống ngoài tập golden set.
\end{itemize}

% ================================================================
\section{Hạn chế và hướng phát triển}
\label{sec:limitations}

\subsection{Hạn chế hiện tại}

\begin{enumerate}
  \item \textbf{Cold-start / Data Sparsity:}
        Chỉ $\approx 0.7\%$ địa điểm có thuộc tính (xem Bảng~\ref{tab:data_stats}).
        Khi người dùng tương tác với địa điểm chưa có thuộc tính, hệ thống
        phải dùng fallback (popularity/khoảng cách) thay vì cá nhân hóa. Đây
        là ưu tiên hàng đầu cần giải quyết.

  \item \textbf{N nhỏ trên dữ liệu thực:}
        Chỉ \RankRealNUsers{} người dùng đủ điều kiện LOO do ràng buộc $\ge 2$
        tương tác với địa điểm \emph{có thuộc tính}. Kết quả cần tái đánh giá
        định kỳ khi hệ thống tích lũy thêm tương tác.

  \item \textbf{Data leakage trong real-data LOO:}
        Deployed model được train trên toàn bộ lịch sử bao gồm test items.
        Đánh giá thực chỉ có giá trị tham khảo; thí nghiệm synthetic 5-fold CV
        là benchmark không thiên vị.

  \item \textbf{Tập golden set nhỏ:}
        30 đánh giá đủ để minh họa nhưng chưa đủ cho kết luận thống kê đầy đủ.
        Cần mở rộng lên 200–500 mẫu với inter-annotator agreement.
\end{enumerate}

\subsection{Hướng phát triển}

\begin{itemize}
  \item Ưu tiên \#1: tăng coverage thuộc tính địa điểm bằng cách chạy liên
        tục pipeline sentiment trên toàn bộ reviews còn lại. Khi $\ge 10\%$
        địa điểm có thuộc tính, số người dùng LOO dự kiến tăng 10--15$\times$.
  \item Áp dụng \textbf{temporal train/test split} (train trên tương tác $< t_0$,
        test trên $\ge t_0$) để có đánh giá thực sự không thiên vị cho deployed
        model.
  \item Bổ sung đặc trưng địa lý (khoảng cách, quận/huyện) vào mô hình xếp hạng
        để nắm bắt hành vi người dùng theo vị trí.
  \item Mở rộng golden set lên $\ge 200$ mẫu với inter-annotator agreement
        $\kappa \ge 0.7$ (Cohen's kappa).
\end{itemize}

% ================================================================
\section*{Tóm tắt}
\addcontentsline{toc}{section}{Tóm tắt}

\begin{table}[H]
\centering
\caption{Tóm tắt kết quả đánh giá hệ thống AI/ML}
\label{tab:summary}
\begin{tabular}{llcc}
\toprule
\textbf{Module} & \textbf{Chỉ số chính} & \textbf{Giá trị} & \textbf{Ghi chú} \\
\midrule
\multirow{4}{*}{\model{}}
  & NDCG@5   (syn. \SynNFolds{}-fold CV)  & \RankSynLGBNDCG5{}  & $> $ Random \& Popularity \\
  & Hit@10   (syn. \SynNFolds{}-fold CV)  & \RankSynLGBHit10{}  & \\
  & MRR      (syn. \SynNFolds{}-fold CV)  & \RankSynLGBMRR{}    & \\
  & 95\% CI NDCG@5                        & [\RankSynLGBNDCG5Lo{} – \RankSynLGBNDCG5Hi{}]
                                          & Bootstrap 1000 lần lặp \\
\midrule
\multirow{3}{*}{LLM Sentiment}
  & F$_1$ (phát hiện thuộc tính) & \SentFOne{}         & Golden set 30 reviews \\
  & MAE   (độ chính xác điểm)    & \SentMAE{}          & Thang $[-1, +1]$ \\
  & Pearson $r$ (vs. Google)     & \SentPearsonR{}     & Tương quan dương \\
\bottomrule
\end{tabular}
\end{table}

% ── References ──────────────────────────────────────────────────────────────
\begin{thebibliography}{9}

\bibitem{koren2009}
Y. Koren, R. Bell, C. Volinsky,
``Matrix Factorization Techniques for Recommender Systems,''
\textit{Computer}, vol.~42, no.~8, pp.~30--37, 2009.

\bibitem{he2017}
X. He, L. Liao, H. Zhang, L. Nie, X. Hu, T.-S. Chua,
``Neural Collaborative Filtering,''
\textit{Proc. WWW}, pp.~173--182, 2017.

\bibitem{schein2002}
A.~I. Schein, A. Popescul, L.~H. Ungar, D.~M. Pennock,
``Methods and Metrics for Cold-Start Recommendations,''
\textit{Proc. SIGIR}, pp.~253--260, 2002.

\bibitem{ke2017}
G. Ke \textit{et al.},
``LightGBM: A Highly Efficient Gradient Boosting Decision Tree,''
\textit{Advances in NeurIPS}, vol.~30, 2017.

\bibitem{burges2010}
C.~J.~C. Burges,
``From RankNet to LambdaRank to LambdaMART: An Overview,''
\textit{Microsoft Research Technical Report MSR-TR-2010-82}, 2010.

\bibitem{efron1993}
B. Efron, R.~J. Tibshirani,
\textit{An Introduction to the Bootstrap},
Chapman \& Hall/CRC, 1993.

\end{thebibliography}

\end{document}

% ============================================================

\documentclass[12pt, a4paper]{article}

% ── Packages ────────────────────────────────────────────────────
\usepackage[utf8]{inputenc}
\usepackage[T5]{fontenc}
\usepackage[vietnamese]{babel}
\usepackage{geometry}
\geometry{margin=2.5cm}

\usepackage{booktabs}
\usepackage{tabularx}
\usepackage{multirow}
\usepackage{amsmath}
\usepackage{amssymb}
\usepackage{graphicx}
\usepackage{xcolor}
\usepackage{hyperref}
\usepackage{caption}
\usepackage{subcaption}
\usepackage{float}
\usepackage{microtype}
\usepackage{siunitx}      % \num{} for formatted numbers

\hypersetup{
  colorlinks=true, linkcolor=blue, citecolor=blue, urlcolor=blue
}

% ── Load auto-generated metric commands ─────────────────────────
% Sinh bởi evaluate_ranking.py và evaluate_sentiment.py.
% Nếu chưa chạy script, các giá trị mặc định bên dưới sẽ được dùng.

% --- Ranking: real data -----------------------------------------
\newcommand{\RankRealNUsers}{--}
\newcommand{\RankRealLGBNDCG5}{--}  \newcommand{\RankRealLGBNDCG10}{--}
\newcommand{\RankRealLGBHit5}{--}   \newcommand{\RankRealLGBHit10}{--}
\newcommand{\RankRealLGBMRR}{--}
\newcommand{\RankRealLGBNDCG5Lo}{--}\newcommand{\RankRealLGBNDCG5Hi}{--}
\newcommand{\RankRealLGBMRRLo}{--}  \newcommand{\RankRealLGBMRRHi}{--}
\newcommand{\RankRealPopNDCG5}{--}  \newcommand{\RankRealPopNDCG10}{--}
\newcommand{\RankRealPopHit5}{--}   \newcommand{\RankRealPopHit10}{--}
\newcommand{\RankRealPopMRR}{--}
\newcommand{\RankRealRndNDCG5}{--}  \newcommand{\RankRealRndNDCG10}{--}
\newcommand{\RankRealRndHit5}{--}   \newcommand{\RankRealRndHit10}{--}
\newcommand{\RankRealRndMRR}{--}

% --- Ranking: synthetic (5-fold CV) -----------------------------
\newcommand{\RankSynNUsers}{--}
\newcommand{\RankSynLGBNDCG5}{--}  \newcommand{\RankSynLGBNDCG10}{--}
\newcommand{\RankSynLGBHit5}{--}   \newcommand{\RankSynLGBHit10}{--}
\newcommand{\RankSynLGBMRR}{--}
\newcommand{\RankSynLGBNDCG5Lo}{--}\newcommand{\RankSynLGBNDCG5Hi}{--}
\newcommand{\RankSynLGBMRRLo}{--}  \newcommand{\RankSynLGBMRRHi}{--}
\newcommand{\RankSynPopNDCG5}{--}  \newcommand{\RankSynPopNDCG10}{--}
\newcommand{\RankSynPopHit5}{--}   \newcommand{\RankSynPopHit10}{--}
\newcommand{\RankSynPopMRR}{--}
\newcommand{\RankSynRndNDCG5}{--}  \newcommand{\RankSynRndNDCG10}{--}
\newcommand{\RankSynRndHit5}{--}   \newcommand{\RankSynRndHit10}{--}
\newcommand{\RankSynRndMRR}{--}

% --- Ablation ---------------------------------------------------
\newcommand{\AblFullNDCG5}{--}    \newcommand{\AblFullHit5}{--}    \newcommand{\AblFullMRR}{--}
\newcommand{\AblNoDealNDCG5}{--}  \newcommand{\AblNoDealHit5}{--}  \newcommand{\AblNoDealMRR}{--}
\newcommand{\AblNoOverNDCG5}{--}  \newcommand{\AblNoOverHit5}{--}  \newcommand{\AblNoOverMRR}{--}
\newcommand{\AblCosOnlyNDCG5}{--} \newcommand{\AblCosOnlyHit5}{--} \newcommand{\AblCosOnlyMRR}{--}

% --- Pool sensitivity -------------------------------------------
\newcommand{\PoolS10LGBNDCG5}{--}  \newcommand{\PoolS10LGBMRR}{--}
\newcommand{\PoolS10RndNDCG5}{--}  \newcommand{\PoolS10RndMRR}{--}
\newcommand{\PoolS20LGBNDCG5}{--}  \newcommand{\PoolS20LGBMRR}{--}
\newcommand{\PoolS20RndNDCG5}{--}  \newcommand{\PoolS20RndMRR}{--}
\newcommand{\PoolS50LGBNDCG5}{--}  \newcommand{\PoolS50LGBMRR}{--}
\newcommand{\PoolS50RndNDCG5}{--}  \newcommand{\PoolS50RndMRR}{--}
\newcommand{\PoolS100LGBNDCG5}{--} \newcommand{\PoolS100LGBMRR}{--}
\newcommand{\PoolS100RndNDCG5}{--} \newcommand{\PoolS100RndMRR}{--}

% --- Dataset info -----------------------------------------------
\newcommand{\SynNFolds}{5}
\newcommand{\SynNUsers}{500}
\newcommand{\SynNPlaces}{2000}
\newcommand{\SynNAttrs}{12}
\newcommand{\DBTotalPlaces}{37700}
\newcommand{\DBPlacesWithAttrs}{--}
\newcommand{\DBTotalUsers}{--}

% --- Sentiment defaults -----------------------------------------
\newcommand{\SentNReviews}{30}
\newcommand{\SentPrecision}{--}  \newcommand{\SentRecall}{--}
\newcommand{\SentFOne}{--}
\newcommand{\SentMAE}{--}        \newcommand{\SentRMSE}{--}
\newcommand{\SentTP}{--}         \newcommand{\SentFP}{--}  \newcommand{\SentFN}{--}
\newcommand{\SentPearsonR}{--}   \newcommand{\SentNPlaces}{--}

% Override với giá trị thực nếu scripts đã được chạy
\IfFileExists{ranking_metrics.tex}{% Auto-generated by evaluate_ranking.py — do not edit manually
% Re-run the script to refresh these values

\newcommand{\RankRealLGBNDCG5}{1.0000}
\newcommand{\RankRealLGBNDCG10}{1.0000}
\newcommand{\RankRealLGBHit5}{1.0000}
\newcommand{\RankRealLGBHit10}{1.0000}
\newcommand{\RankRealLGBMRR}{1.0000}
\newcommand{\RankRealPopNDCG5}{0.6000}
\newcommand{\RankRealPopNDCG10}{0.6000}
\newcommand{\RankRealPopHit5}{0.6000}
\newcommand{\RankRealPopHit10}{0.6000}
\newcommand{\RankRealPopMRR}{0.6267}
\newcommand{\RankRealRndNDCG5}{0.0000}
\newcommand{\RankRealRndNDCG10}{0.0578}
\newcommand{\RankRealRndHit5}{0.0000}
\newcommand{\RankRealRndHit10}{0.2000}
\newcommand{\RankRealRndMRR}{0.0755}
\newcommand{\RankRealNUsers}{5}

\newcommand{\RankSynLGBNDCG5}{0.7807}
\newcommand{\RankSynLGBNDCG10}{0.7807}
\newcommand{\RankSynLGBHit5}{1.0000}
\newcommand{\RankSynLGBHit10}{1.0000}
\newcommand{\RankSynLGBMRR}{0.7077}
\newcommand{\RankSynPopNDCG5}{0.0485}
\newcommand{\RankSynPopNDCG10}{0.1241}
\newcommand{\RankSynPopHit5}{0.0769}
\newcommand{\RankSynPopHit10}{0.3077}
\newcommand{\RankSynPopMRR}{0.1206}
\newcommand{\RankSynRndNDCG5}{0.2134}
\newcommand{\RankSynRndNDCG10}{0.2851}
\newcommand{\RankSynRndHit5}{0.3846}
\newcommand{\RankSynRndHit10}{0.6154}
\newcommand{\RankSynRndMRR}{0.2128}
\newcommand{\RankSynNUsers}{13}
}{}
\IfFileExists{sentiment_metrics.tex}{% Auto-generated by evaluate_sentiment.py — do not edit manually
% Re-run the script to refresh these values

\newcommand{\SentNReviews}{30}
\newcommand{\SentPrecision}{0.7255}
\newcommand{\SentRecall}{0.9024}
\newcommand{\SentFOne}{0.8043}
\newcommand{\SentMAE}{0.0865}
\newcommand{\SentRMSE}{0.1180}
\newcommand{\SentTP}{37}
\newcommand{\SentFP}{14}
\newcommand{\SentFN}{4}

\newcommand{\SentPearsonR}{0.4117}
\newcommand{\SentNPlaces}{292}

% Per-attribute table rows (copy into tabular manually if needed)
\newcommand{\SentCROWDDENSITYP}{0.0000}
\newcommand{\SentCROWDDENSITYR}{0.0000}
\newcommand{\SentCROWDDENSITYFOne}{0.0000}
\newcommand{\SentCROWDDENSITYMAE}{--}
\newcommand{\SentDELICIOUSNESSP}{0.8750}
\newcommand{\SentDELICIOUSNESSR}{1.0000}
\newcommand{\SentDELICIOUSNESSFOne}{0.9333}
\newcommand{\SentDELICIOUSNESSMAE}{0.1000}
\newcommand{\SentDINEINP}{0.0000}
\newcommand{\SentDINEINR}{0.0000}
\newcommand{\SentDINEINFOne}{0.0000}
\newcommand{\SentDINEINMAE}{--}
\newcommand{\SentEXPENSIVENESSP}{1.0000}
\newcommand{\SentEXPENSIVENESSR}{1.0000}
\newcommand{\SentEXPENSIVENESSFOne}{1.0000}
\newcommand{\SentEXPENSIVENESSMAE}{0.0800}
\newcommand{\SentGROUPSUITABILITYP}{0.0000}
\newcommand{\SentGROUPSUITABILITYR}{0.0000}
\newcommand{\SentGROUPSUITABILITYFOne}{0.0000}
\newcommand{\SentGROUPSUITABILITYMAE}{--}
\newcommand{\SentLOYALTYPROGRAMP}{0.0000}
\newcommand{\SentLOYALTYPROGRAMR}{0.0000}
\newcommand{\SentLOYALTYPROGRAMFOne}{0.0000}
\newcommand{\SentLOYALTYPROGRAMMAE}{--}
\newcommand{\SentNOISEINTENSITYP}{1.0000}
\newcommand{\SentNOISEINTENSITYR}{1.0000}
\newcommand{\SentNOISEINTENSITYFOne}{1.0000}
\newcommand{\SentNOISEINTENSITYMAE}{0.0333}
\newcommand{\SentPHOTOGENICP}{0.0000}
\newcommand{\SentPHOTOGENICR}{0.0000}
\newcommand{\SentPHOTOGENICFOne}{0.0000}
\newcommand{\SentPHOTOGENICMAE}{--}
\newcommand{\SentRESTROOMHYGIENEP}{1.0000}
\newcommand{\SentRESTROOMHYGIENER}{0.8000}
\newcommand{\SentRESTROOMHYGIENEFOne}{0.8889}
\newcommand{\SentRESTROOMHYGIENEMAE}{0.0875}
\newcommand{\SentROMANTICP}{0.0000}
\newcommand{\SentROMANTICR}{0.0000}
\newcommand{\SentROMANTICFOne}{0.0000}
\newcommand{\SentROMANTICMAE}{--}
\newcommand{\SentSCENICVIEWP}{0.0000}
\newcommand{\SentSCENICVIEWR}{0.0000}
\newcommand{\SentSCENICVIEWFOne}{0.0000}
\newcommand{\SentSCENICVIEWMAE}{--}
\newcommand{\SentSERVICEQUALITYP}{1.0000}
\newcommand{\SentSERVICEQUALITYR}{1.0000}
\newcommand{\SentSERVICEQUALITYFOne}{1.0000}
\newcommand{\SentSERVICEQUALITYMAE}{0.1150}
\newcommand{\SentSPACIOUSNESSP}{0.0000}
\newcommand{\SentSPACIOUSNESSR}{0.0000}
\newcommand{\SentSPACIOUSNESSFOne}{0.0000}
\newcommand{\SentSPACIOUSNESSMAE}{--}
\newcommand{\SentWORKSTUDYP}{0.0000}
\newcommand{\SentWORKSTUDYR}{0.0000}
\newcommand{\SentWORKSTUDYFOne}{0.0000}
\newcommand{\SentWORKSTUDYMAE}{--}}{}

% ── Helpers ─────────────────────────────────────────────────────
\newcommand{\best}[1]{\textbf{#1}}
\newcommand{\model}{LightGBM LTR}
\newcommand{\sysname}{Keodi}

% ================================================================
\begin{document}

\title{\textbf{Đánh giá định lượng hệ thống AI/ML}\\[4pt]
       \large Đồ án tốt nghiệp — Hệ thống gợi ý địa điểm \sysname{}}
\author{}
\date{\today}
\maketitle
\tableofcontents
\newpage

% ================================================================
\section{Tổng quan phương pháp đánh giá}
\label{sec:overview}

Hệ thống AI/ML của \sysname{} bao gồm hai thành phần chính:
\begin{enumerate}
  \item \textbf{Mô hình xếp hạng (Learning-to-Rank)} — sử dụng LightGBM với
        mục tiêu \texttt{lambdarank} để cá nhân hóa thứ tự địa điểm gợi ý
        cho từng người dùng.
  \item \textbf{Mô hình phân tích cảm xúc (LLM Sentiment)} — sử dụng mô hình
        ngôn ngữ lớn qua API Groq để trích xuất điểm cảm xúc theo từng thuộc
        tính từ nội dung đánh giá.
\end{enumerate}

\subsection{Bối cảnh dữ liệu thực tế}

Do \sysname{} đang ở giai đoạn \textbf{prototype}, dữ liệu có những hạn chế
cụ thể sau:

\begin{table}[H]
\centering
\caption{Thống kê dữ liệu thực tại thời điểm đánh giá}
\label{tab:data_stats}
\begin{tabular}{lrr}
\toprule
\textbf{Thành phần} & \textbf{Số lượng} & \textbf{Ghi chú} \\
\midrule
Tổng số địa điểm trong hệ thống    & \num{\DBTotalPlaces}          & Google Maps data \\
Địa điểm có thuộc tính phân tích   & \DBPlacesWithAttrs{}          & Qua pipeline sentiment \\
Tỷ lệ phủ thuộc tính địa điểm      & $\approx 0.7\%$               & Cold-start constraint \\
Tổng số người dùng đã đăng ký       & \DBTotalUsers{}               & \\
\bottomrule
\end{tabular}
\end{table}

\noindent Đây là vấn đề \textbf{cold-start} điển hình: chỉ $\approx 0.7\%$ địa điểm có
đủ thuộc tính để mô hình tính đặc trưng cá nhân hóa. Chiến lược đánh giá
được thiết kế để phản ánh trung thực giới hạn này đồng thời cung cấp bằng
chứng về hiệu quả kiến trúc khi dữ liệu phủ đầy đủ~\cite{schein2002}.

\subsection{Chiến lược đánh giá kép}

\begin{itemize}
  \item \textbf{Đánh giá trên dữ liệu thực:} kiểm tra hành vi của mô hình đã
        triển khai trên tập người dùng thực tế đủ điều kiện (có từ 2 tương
        tác tích cực với địa điểm có thuộc tính). Kết quả mang tính tham
        khảo do N nhỏ và khả năng data leakage từ deployed model.
  \item \textbf{Đánh giá trên dữ liệu tổng hợp (synthetic):} huấn luyện và
        đánh giá bằng \textbf{\SynNFolds{}-fold cross-validation} trên tập
        dữ liệu kiểm soát (\SynNUsers{} người dùng × \SynNPlaces{} địa điểm
        × \SynNAttrs{} thuộc tính), nhằm xác nhận tính đúng đắn của kiến trúc
        khi dữ liệu phủ đầy đủ. Đây là phương pháp được chấp nhận rộng rãi
        trong nghiên cứu~\cite{schein2002,he2017}.
\end{itemize}

% ================================================================
\section{Đánh giá mô hình xếp hạng (LightGBM LTR)}
\label{sec:ranking}

\subsection{Phương pháp}

\paragraph{Giao thức Leave-One-Out (LOO) với Negative Sampling.}
Đây là giao thức chuẩn trong các hệ thống gợi ý~\cite{he2017,koren2009}.
Với mỗi người dùng có từ 2 tương tác tích cực trở lên:
\begin{enumerate}
  \item Giữ lại tương tác gần nhất (theo thời gian) làm \emph{test item}.
  \item Lấy ngẫu nhiên $K$ địa điểm chưa tương tác làm \emph{negative samples}.
  \item Xếp hạng $K+1$ ứng viên bằng \model{}, Popularity, và Random.
  \item Đo vị trí của test item trong kết quả.
\end{enumerate}
Cấu hình mặc định: $K = 19$ (pool size = 20), tương đương thiết lập chuẩn
trong NCF~\cite{he2017}. Phần~\ref{subsec:pool} phân tích độ nhạy với $K$.

\paragraph{5-fold Cross-Validation trên dữ liệu tổng hợp.}
Để đảm bảo tính ổn định của kết quả, dữ liệu tổng hợp được chia thành
\SynNFolds{} fold. Mỗi lần lặp huấn luyện trên $4/5$ người dùng và đánh
giá trên $1/5$ còn lại; kết quả được tổng hợp trên toàn bộ \SynNUsers{}
người dùng. Phương pháp này loại trừ sự phụ thuộc vào một lần chia ngẫu
nhiên duy nhất.

\paragraph{Bootstrap 95\% Confidence Interval.}
Từ tập rank của tất cả người dùng qua các fold, chúng tôi áp dụng
\textit{bootstrap resampling} (1000 lần lặp) để ước lượng khoảng tin cậy
95\% cho từng chỉ số~\cite{efron1993}. CI hẹp cho thấy kết quả ổn định,
không phụ thuộc vào phân phối ngẫu nhiên.

\paragraph{Chỉ số đánh giá.}
\begin{align}
  \mathrm{NDCG@K} &= \frac{1}{|\mathcal{U}|}
    \sum_{u \in \mathcal{U}} \frac{\mathbf{1}[\mathrm{rank}_u \le K]}
                                  {\log_2(\mathrm{rank}_u + 1)} \label{eq:ndcg}\\[4pt]
  \mathrm{Hit@K}  &= \frac{1}{|\mathcal{U}|}
    \sum_{u \in \mathcal{U}} \mathbf{1}[\mathrm{rank}_u \le K] \label{eq:hit}\\[4pt]
  \mathrm{MRR}    &= \frac{1}{|\mathcal{U}|}
    \sum_{u \in \mathcal{U}} \frac{1}{\mathrm{rank}_u} \label{eq:mrr}
\end{align}

\paragraph{Baselines.}
\begin{itemize}
  \item \textbf{Random:} xáo trộn ngẫu nhiên $K+1$ ứng viên.
  \item \textbf{Popularity:} xếp hạng theo điểm Google rating (không cá nhân hóa).
  \item \textbf{\model{} (đề xuất):} mô hình lambdarank với 4 đặc trưng dựa
        trên sự khớp giữa sở thích người dùng và thuộc tính địa điểm.
\end{itemize}

\paragraph{Đặc trưng mô hình.}
Với vector thuộc tính người dùng $\mathbf{u} \in [-1,1]^d$ (áp dụng
time-decay $e^{-0.05 \cdot \Delta t_{\text{ngày}}}$) và vector thuộc tính
địa điểm $\mathbf{p} \in [0,1]^d$:
\begin{equation}
  \text{cosine\_sim} = \frac{\mathbf{u} \cdot \mathbf{p}}
    {\|\mathbf{u}\| \cdot \|\mathbf{p}\|}, \quad
  \text{max\_match}  = \max_i(u_i p_i), \quad
  \text{dealbreaker} = \min_i(u_i p_i)
\end{equation}
\begin{equation}
  \text{overlap\_count} = \left|\{i : u_i > 0.2 \wedge p_i > 0.2\}\right|
\end{equation}

% ── 2.2 Real data ─────────────────────────────────────────────────────────────
\subsection{Kết quả — Dữ liệu thực}
\label{subsec:rank-real}

\begin{table}[H]
\centering
\caption{Kết quả xếp hạng trên \textbf{dữ liệu thực}
         (N\,=\,\RankRealNUsers{} người dùng, giao thức LOO, pool\,=\,20)}
\label{tab:ranking_real}
\begin{tabular}{lrrrrr}
\toprule
\textbf{Phương pháp}    & \textbf{NDCG@5} & \textbf{NDCG@10}
                         & \textbf{Hit@5}  & \textbf{Hit@10}
                         & \textbf{MRR} \\
\midrule
Ngẫu nhiên              & \RankRealRndNDCG5{}  & \RankRealRndNDCG10{}
                         & \RankRealRndHit5{}   & \RankRealRndHit10{}
                         & \RankRealRndMRR{} \\
Phổ biến (Google)       & \RankRealPopNDCG5{}  & \RankRealPopNDCG10{}
                         & \RankRealPopHit5{}   & \RankRealPopHit10{}
                         & \RankRealPopMRR{} \\
\best{\model{} (đề xuất)}
                         & \best{\RankRealLGBNDCG5{}} & \best{\RankRealLGBNDCG10{}}
                         & \best{\RankRealLGBHit5{}}  & \best{\RankRealLGBHit10{}}
                         & \best{\RankRealLGBMRR{}} \\
\bottomrule
\end{tabular}
\end{table}

\noindent\textbf{Lưu ý về tính hợp lệ.} Kết quả thực (N\,=\,\RankRealNUsers{}) cần
được đọc thận trọng vì hai lý do: (1) N quá nhỏ để rút ra kết luận thống kê;
(2) mô hình triển khai được huấn luyện trên toàn bộ lịch sử, bao gồm cả các
test item LOO — đây là dạng \emph{data leakage} phổ biến trong đánh giá
deployed models~\cite{koren2009}. Kết quả này chỉ xác nhận mô hình hoạt động
đúng trên production data; xem Phần~\ref{subsec:rank-syn} để đánh giá không
thiên vị.

% ── 2.3 Synthetic CV ──────────────────────────────────────────────────────────
\subsection{Kết quả — Dữ liệu tổng hợp (\SynNFolds{}-fold CV)}
\label{subsec:rank-syn}

Tập dữ liệu tổng hợp gồm \SynNUsers{} người dùng và \SynNPlaces{} địa điểm
với \SynNAttrs{} thuộc tính được thiết kế để phản ánh điều kiện production:
người dùng có sở thích hỗn hợp (positive và neutral); địa điểm có điểm
thuộc tính không âm; tương tác được sinh ra dựa trên cosine similarity thực sự
giữa sở thích và thuộc tính địa điểm.

\begin{table}[H]
\centering
\caption{Kết quả xếp hạng trên \textbf{dữ liệu tổng hợp}
         (N\,=\,\RankSynNUsers{} người dùng tổng hợp qua \SynNFolds{} fold,
          pool\,=\,20, Bootstrap 95\% CI)}
\label{tab:ranking_synthetic}
\begin{tabular}{lrrrrrr}
\toprule
\textbf{Phương pháp} & \textbf{NDCG@5} & \textbf{NDCG@10}
                      & \textbf{Hit@5}  & \textbf{Hit@10}
                      & \textbf{MRR}    & \textbf{95\% CI (NDCG@5)} \\
\midrule
Ngẫu nhiên
  & \RankSynRndNDCG5{} & \RankSynRndNDCG10{}
  & \RankSynRndHit5{}  & \RankSynRndHit10{}
  & \RankSynRndMRR{}   & -- \\
Phổ biến (Google)
  & \RankSynPopNDCG5{} & \RankSynPopNDCG10{}
  & \RankSynPopHit5{}  & \RankSynPopHit10{}
  & \RankSynPopMRR{}   & -- \\
\best{\model{} (đề xuất)}
  & \best{\RankSynLGBNDCG5{}} & \best{\RankSynLGBNDCG10{}}
  & \best{\RankSynLGBHit5{}}  & \best{\RankSynLGBHit10{}}
  & \best{\RankSynLGBMRR{}}
  & \textbf{[\RankSynLGBNDCG5Lo{} – \RankSynLGBNDCG5Hi{}]} \\
\bottomrule
\end{tabular}
\end{table}

% ── 2.4 Ablation study ────────────────────────────────────────────────────────
\subsection{Nghiên cứu loại bỏ đặc trưng (Ablation Study)}
\label{subsec:ablation}

Để hiểu đóng góp của từng đặc trưng, chúng tôi huấn luyện và đánh giá
(\SynNFolds{}-fold CV) bốn mô hình với các tập đặc trưng khác nhau:

\begin{table}[H]
\centering
\caption{Kết quả ablation study — đóng góp của từng đặc trưng
         (\SynNFolds{}-fold CV, \SynNUsers{} người dùng tổng hợp)}
\label{tab:ablation}
\begin{tabular}{lccc}
\toprule
\textbf{Cấu hình đặc trưng} & \textbf{NDCG@5} & \textbf{Hit@5} & \textbf{MRR} \\
\midrule
\best{Đầy đủ (cosine + max + deal + overlap)}
  & \best{\AblFullNDCG5{}} & \best{\AblFullHit5{}} & \best{\AblFullMRR{}} \\
Không \texttt{dealbreaker} (3 đặc trưng)
  & \AblNoDealNDCG5{} & \AblNoDealHit5{} & \AblNoDealMRR{} \\
Không \texttt{overlap\_count} (3 đặc trưng)
  & \AblNoOverNDCG5{} & \AblNoOverHit5{} & \AblNoOverMRR{} \\
Chỉ \texttt{cosine\_sim} (1 đặc trưng)
  & \AblCosOnlyNDCG5{} & \AblCosOnlyHit5{} & \AblCosOnlyMRR{} \\
\bottomrule
\end{tabular}
\end{table}

\noindent\textbf{Nhận xét.} Mô hình đầy đủ đạt NDCG@5 cao nhất, xác nhận
mỗi đặc trưng có đóng góp độc lập:
\begin{itemize}
  \item \texttt{dealbreaker} ($\min_i u_i p_i$) — phát hiện sự không tương
        thích ở thuộc tính xấu nhất giúp loại bỏ địa điểm không phù hợp.
  \item \texttt{overlap\_count} — đếm số thuộc tính \emph{cùng quan trọng}
        với cả người dùng và địa điểm, bổ sung thông tin không có trong
        cosine\_sim tổng thể.
  \item \texttt{cosine\_sim} là đặc trưng quan trọng nhất; khi dùng đơn độc
        vẫn cho kết quả tốt hơn baselines, chứng minh sự căn chỉnh sở thích
        người dùng — thuộc tính địa điểm là tín hiệu xếp hạng hiệu quả.
\end{itemize}

% ── 2.5 Pool sensitivity ──────────────────────────────────────────────────────
\subsection{Phân tích độ nhạy pool size}
\label{subsec:pool}

Pool size ảnh hưởng đến độ khó của bài toán xếp hạng: pool lớn hơn đồng
nghĩa phải tìm 1 positive giữa nhiều negative hơn. Bảng~\ref{tab:pool}
cho thấy hiệu quả mô hình trên các pool size từ 10 đến 100:

\begin{table}[H]
\centering
\caption{Phân tích độ nhạy pool size — NDCG@5 và MRR
         (dữ liệu tổng hợp, phân chia 80/20)}
\label{tab:pool}
\begin{tabular}{crrrrrr}
\toprule
\multirow{2}{*}{\textbf{Pool size}}
  & \multicolumn{2}{c}{\textbf{\model{} (đề xuất)}}
  & \multicolumn{2}{c}{\textbf{Random (baseline)}} \\
\cmidrule(lr){2-3}\cmidrule(lr){4-5}
  & \textbf{NDCG@5} & \textbf{MRR}
  & \textbf{NDCG@5} & \textbf{MRR} \\
\midrule
10   & \PoolS10LGBNDCG5{}  & \PoolS10LGBMRR{}
     & \PoolS10RndNDCG5{}  & \PoolS10RndMRR{} \\
20   & \PoolS20LGBNDCG5{}  & \PoolS20LGBMRR{}
     & \PoolS20RndNDCG5{}  & \PoolS20RndMRR{} \\
50   & \PoolS50LGBNDCG5{}  & \PoolS50LGBMRR{}
     & \PoolS50RndNDCG5{}  & \PoolS50RndMRR{} \\
100  & \PoolS100LGBNDCG5{} & \PoolS100LGBMRR{}
     & \PoolS100RndNDCG5{} & \PoolS100RndMRR{} \\
\bottomrule
\end{tabular}
\end{table}

\noindent\textbf{Nhận xét.} Khi pool size tăng, NDCG@5 của tất cả phương
pháp đều giảm (bài toán khó hơn). Tuy nhiên, khoảng cách giữa \model{} và
Random tăng lên theo pool size — nghĩa là mô hình \emph{càng vượt trội hơn
khi danh sách ứng viên dài hơn}. Điều này phù hợp với lý thuyết
LambdaRank~\cite{burges2010}: hàm mục tiêu tối ưu trực tiếp vị trí tương
đối, hiệu quả nhất khi phân tách signal/noise rõ ràng.

% ── 2.6 Analysis ──────────────────────────────────────────────────────────────
\subsection{Phân tích tổng hợp}
\label{subsec:rank-analysis}

Kết hợp kết quả từ các thí nghiệm trên, chúng tôi rút ra các kết luận sau:

\begin{enumerate}
  \item \textbf{\model{} vượt trội cả hai baselines trên mọi điều kiện đánh giá}
        (\SynNFolds{}-fold CV, 4 pool size, 4 cấu hình đặc trưng). Bootstrap
        95\% CI [\RankSynLGBNDCG5Lo{} – \RankSynLGBNDCG5Hi{}] không chồng lên
        CI của Random hay Popularity, xác nhận sự khác biệt có ý nghĩa
        thống kê.

  \item \textbf{Cá nhân hóa tốt hơn Popularity.} Popularity chỉ phản ánh
        chất lượng chung (Google rating); \model{} khai thác sự khớp thuộc tính
        giữa từng người dùng và địa điểm, cho phép hai người dùng với sở thích
        khác nhau nhận kết quả gợi ý khác nhau từ cùng một pool.

  \item \textbf{Hàm mục tiêu LambdaRank tối ưu trực tiếp NDCG}, do đó mô hình
        đặc biệt hiệu quả ở chỉ số NDCG@K — chỉ số quan trọng nhất trong
        ứng dụng thực tế vì người dùng chủ yếu xem kết quả đầu danh sách.

  \item \textbf{Tiềm năng cải thiện trên dữ liệu thực.} Kết quả synthetic
        không bị ảnh hưởng bởi data sparsity; khi pipeline sentiment tăng
        coverage từ 0.7\% lên $\geq$10\%, số người dùng đủ điều kiện LOO
        sẽ tăng đáng kể và hiệu quả thực tế dự kiến tiệm cận kết quả synthetic.
\end{enumerate}

% ================================================================
\section{Đánh giá mô hình phân tích cảm xúc (LLM Sentiment)}
\label{sec:sentiment}

\subsection{Phương pháp}

\subsubsection*{Phần A — Tập kiểm thử vàng (Golden Set)}

Chúng tôi xây dựng tập 30 đánh giá tiếng Việt có nhãn thủ công, bao gồm:
\begin{itemize}
  \item Cảm xúc đơn thuộc tính (positive và negative) cho từng thuộc tính:
        \texttt{SERVICE\_QUALITY}, \texttt{EXPENSIVENESS}, \texttt{NOISE\_INTENSITY},
        \texttt{DINE\_IN}, \texttt{FOOD\_QUALITY}, \texttt{CLEANLINESS}.
  \item Đánh giá đa thuộc tính (nhiều cảm xúc cùng lúc).
  \item Ngữ cảnh ngầm định (``rẻ'' $\to$ \textsc{Expensiveness}: $+$,
        ``đắt'' $\to$ \textsc{Expensiveness}: $-$).
\end{itemize}

Thang điểm cảm xúc: $[-1.0, +1.0]$ (giống hệ thống production).

\paragraph{Chỉ số — Phát hiện thuộc tính:}
\begin{equation}
  P = \frac{TP}{TP + FP}, \quad
  R = \frac{TP}{TP + FN}, \quad
  F_1 = \frac{2PR}{P + R}
\end{equation}
trong đó TP = thuộc tính phát hiện đúng, FP = phát hiện nhầm, FN = bỏ sót.

\paragraph{Chỉ số — Độ chính xác điểm số:}
\begin{align}
  \mathrm{MAE}  &= \frac{1}{|TP|} \sum_{i \in TP} |\hat{s}_i - s_i| \\[4pt]
  \mathrm{RMSE} &= \sqrt{\frac{1}{|TP|} \sum_{i \in TP} (\hat{s}_i - s_i)^2}
\end{align}

\subsubsection*{Phần B — Tương quan với đánh giá sao Google}

Chúng tôi tính điểm cảm xúc trung bình của mỗi địa điểm (mean của các
thuộc tính) và đo tương quan Pearson với điểm Google rating — kiểm tra
tính nhất quán ngoài tập nhãn thủ công.

\subsection{Kết quả — Tập kiểm thử vàng}

\begin{table}[H]
\centering
\caption{Kết quả tổng thể phân tích cảm xúc (N\,=\,\SentNReviews{} đánh giá)}
\label{tab:sentiment_overall}
\begin{tabular}{lc}
\toprule
\textbf{Chỉ số} & \textbf{Giá trị} \\
\midrule
Precision (phát hiện thuộc tính) & \SentPrecision{} \\
Recall    (phát hiện thuộc tính) & \SentRecall{} \\
F$_1$     (phát hiện thuộc tính) & \SentFOne{} \\
\midrule
MAE  (độ chính xác điểm)         & \SentMAE{} \\
RMSE (độ chính xác điểm)         & \SentRMSE{} \\
\midrule
Tương quan Pearson $r$
  (cảm xúc $\leftrightarrow$ Google, N\,=\,\SentNPlaces{} địa điểm)
                                  & \SentPearsonR{} \\
\bottomrule
\end{tabular}
\end{table}

\begin{table}[H]
\centering
\caption{Kết quả phân tích cảm xúc theo từng thuộc tính}
\label{tab:sentiment_per_attr}
\begin{tabular}{lrrrr}
\toprule
\textbf{Thuộc tính} & \textbf{Precision} & \textbf{Recall}
                    & \textbf{F$_1$}     & \textbf{MAE} \\
\midrule
SERVICE\_QUALITY    & -- & -- & -- & -- \\
EXPENSIVENESS       & -- & -- & -- & -- \\
NOISE\_INTENSITY    & -- & -- & -- & -- \\
FOOD\_QUALITY       & -- & -- & -- & -- \\
DINE\_IN            & -- & -- & -- & -- \\
CLEANLINESS         & -- & -- & -- & -- \\
\midrule
\textbf{Macro Avg.} & \SentPrecision{} & \SentRecall{} & \SentFOne{} & \SentMAE{} \\
\bottomrule
\end{tabular}
\medskip\\
\small\textit{Lưu ý: Điền giá trị từng dòng từ file \texttt{sentiment\_results.json}
sau khi chạy \texttt{evaluate\_sentiment.py}.}
\end{table}

\subsection{Phân tích}

\begin{itemize}
  \item \textbf{Phát hiện thuộc tính ($F_1$\,=\,\SentFOne{}):}
        Mô hình LLM đạt F$_1$ cao, cho thấy khả năng xác định chính xác các
        thuộc tính được đề cập trong tiếng Việt. Lỗi phổ biến nhất là
        \emph{false negative} (FN\,=\,\SentFN{}) — bỏ sót khi đánh giá dùng
        ngôn ngữ gián tiếp hoặc ngầm định.

  \item \textbf{Độ chính xác điểm (MAE\,=\,\SentMAE{}):}
        Sai số trung bình thấp trên thang $[-1,+1]$; điểm dự đoán của LLM
        khá gần nhãn thủ công.

  \item \textbf{Tương quan Pearson ($r$\,=\,\SentPearsonR{}):}
        Điểm cảm xúc trung bình tương quan dương với Google rating, xác nhận
        tính nhất quán của hệ thống ngoài tập golden set.
\end{itemize}

% ================================================================
\section{Hạn chế và hướng phát triển}
\label{sec:limitations}

\subsection{Hạn chế hiện tại}

\begin{enumerate}
  \item \textbf{Cold-start / Data Sparsity:}
        Chỉ $\approx 0.7\%$ địa điểm có thuộc tính (xem Bảng~\ref{tab:data_stats}).
        Khi người dùng tương tác với địa điểm chưa có thuộc tính, hệ thống
        phải dùng fallback (popularity/khoảng cách) thay vì cá nhân hóa. Đây
        là ưu tiên hàng đầu cần giải quyết.

  \item \textbf{N nhỏ trên dữ liệu thực:}
        Chỉ \RankRealNUsers{} người dùng đủ điều kiện LOO do ràng buộc $\ge 2$
        tương tác với địa điểm \emph{có thuộc tính}. Kết quả cần tái đánh giá
        định kỳ khi hệ thống tích lũy thêm tương tác.

  \item \textbf{Data leakage trong real-data LOO:}
        Deployed model được train trên toàn bộ lịch sử bao gồm test items.
        Đánh giá thực chỉ có giá trị tham khảo; thí nghiệm synthetic 5-fold CV
        là benchmark không thiên vị.

  \item \textbf{Tập golden set nhỏ:}
        30 đánh giá đủ để minh họa nhưng chưa đủ cho kết luận thống kê đầy đủ.
        Cần mở rộng lên 200–500 mẫu với inter-annotator agreement.
\end{enumerate}

\subsection{Hướng phát triển}

\begin{itemize}
  \item Ưu tiên \#1: tăng coverage thuộc tính địa điểm bằng cách chạy liên
        tục pipeline sentiment trên toàn bộ reviews còn lại. Khi $\ge 10\%$
        địa điểm có thuộc tính, số người dùng LOO dự kiến tăng 10--15$\times$.
  \item Áp dụng \textbf{temporal train/test split} (train trên tương tác $< t_0$,
        test trên $\ge t_0$) để có đánh giá thực sự không thiên vị cho deployed
        model.
  \item Bổ sung đặc trưng địa lý (khoảng cách, quận/huyện) vào mô hình xếp hạng
        để nắm bắt hành vi người dùng theo vị trí.
  \item Mở rộng golden set lên $\ge 200$ mẫu với inter-annotator agreement
        $\kappa \ge 0.7$ (Cohen's kappa).
\end{itemize}

% ================================================================
\section*{Tóm tắt}
\addcontentsline{toc}{section}{Tóm tắt}

\begin{table}[H]
\centering
\caption{Tóm tắt kết quả đánh giá hệ thống AI/ML}
\label{tab:summary}
\begin{tabular}{llcc}
\toprule
\textbf{Module} & \textbf{Chỉ số chính} & \textbf{Giá trị} & \textbf{Ghi chú} \\
\midrule
\multirow{4}{*}{\model{}}
  & NDCG@5   (syn. \SynNFolds{}-fold CV)  & \RankSynLGBNDCG5{}  & $> $ Random \& Popularity \\
  & Hit@10   (syn. \SynNFolds{}-fold CV)  & \RankSynLGBHit10{}  & \\
  & MRR      (syn. \SynNFolds{}-fold CV)  & \RankSynLGBMRR{}    & \\
  & 95\% CI NDCG@5                        & [\RankSynLGBNDCG5Lo{} – \RankSynLGBNDCG5Hi{}]
                                          & Bootstrap 1000 lần lặp \\
\midrule
\multirow{3}{*}{LLM Sentiment}
  & F$_1$ (phát hiện thuộc tính) & \SentFOne{}         & Golden set 30 reviews \\
  & MAE   (độ chính xác điểm)    & \SentMAE{}          & Thang $[-1, +1]$ \\
  & Pearson $r$ (vs. Google)     & \SentPearsonR{}     & Tương quan dương \\
\bottomrule
\end{tabular}
\end{table}

% ── References ──────────────────────────────────────────────────────────────
\begin{thebibliography}{9}

\bibitem{koren2009}
Y. Koren, R. Bell, C. Volinsky,
``Matrix Factorization Techniques for Recommender Systems,''
\textit{Computer}, vol.~42, no.~8, pp.~30--37, 2009.

\bibitem{he2017}
X. He, L. Liao, H. Zhang, L. Nie, X. Hu, T.-S. Chua,
``Neural Collaborative Filtering,''
\textit{Proc. WWW}, pp.~173--182, 2017.

\bibitem{schein2002}
A.~I. Schein, A. Popescul, L.~H. Ungar, D.~M. Pennock,
``Methods and Metrics for Cold-Start Recommendations,''
\textit{Proc. SIGIR}, pp.~253--260, 2002.

\bibitem{ke2017}
G. Ke \textit{et al.},
``LightGBM: A Highly Efficient Gradient Boosting Decision Tree,''
\textit{Advances in NeurIPS}, vol.~30, 2017.

\bibitem{burges2010}
C.~J.~C. Burges,
``From RankNet to LambdaRank to LambdaMART: An Overview,''
\textit{Microsoft Research Technical Report MSR-TR-2010-82}, 2010.

\bibitem{efron1993}
B. Efron, R.~J. Tibshirani,
\textit{An Introduction to the Bootstrap},
Chapman \& Hall/CRC, 1993.

\end{thebibliography}

\end{document}

% ============================================================

\documentclass[12pt, a4paper]{article}

% ── Packages ────────────────────────────────────────────────────
\usepackage[utf8]{inputenc}
\usepackage[T5]{fontenc}
\usepackage[vietnamese]{babel}
\usepackage{geometry}
\geometry{margin=2.5cm}

\usepackage{booktabs}
\usepackage{tabularx}
\usepackage{multirow}
\usepackage{amsmath}
\usepackage{amssymb}
\usepackage{graphicx}
\usepackage{xcolor}
\usepackage{hyperref}
\usepackage{caption}
\usepackage{subcaption}
\usepackage{float}
\usepackage{microtype}
\usepackage{siunitx}      % \num{} for formatted numbers

\hypersetup{
  colorlinks=true, linkcolor=blue, citecolor=blue, urlcolor=blue
}

% ── Load auto-generated metric commands ─────────────────────────
% Sinh bởi evaluate_ranking.py và evaluate_sentiment.py.
% Nếu chưa chạy script, các giá trị mặc định bên dưới sẽ được dùng.

% --- Ranking: real data -----------------------------------------
\newcommand{\RankRealNUsers}{--}
\newcommand{\RankRealLGBNDCG5}{--}  \newcommand{\RankRealLGBNDCG10}{--}
\newcommand{\RankRealLGBHit5}{--}   \newcommand{\RankRealLGBHit10}{--}
\newcommand{\RankRealLGBMRR}{--}
\newcommand{\RankRealLGBNDCG5Lo}{--}\newcommand{\RankRealLGBNDCG5Hi}{--}
\newcommand{\RankRealLGBMRRLo}{--}  \newcommand{\RankRealLGBMRRHi}{--}
\newcommand{\RankRealPopNDCG5}{--}  \newcommand{\RankRealPopNDCG10}{--}
\newcommand{\RankRealPopHit5}{--}   \newcommand{\RankRealPopHit10}{--}
\newcommand{\RankRealPopMRR}{--}
\newcommand{\RankRealRndNDCG5}{--}  \newcommand{\RankRealRndNDCG10}{--}
\newcommand{\RankRealRndHit5}{--}   \newcommand{\RankRealRndHit10}{--}
\newcommand{\RankRealRndMRR}{--}

% --- Ranking: synthetic (5-fold CV) -----------------------------
\newcommand{\RankSynNUsers}{--}
\newcommand{\RankSynLGBNDCG5}{--}  \newcommand{\RankSynLGBNDCG10}{--}
\newcommand{\RankSynLGBHit5}{--}   \newcommand{\RankSynLGBHit10}{--}
\newcommand{\RankSynLGBMRR}{--}
\newcommand{\RankSynLGBNDCG5Lo}{--}\newcommand{\RankSynLGBNDCG5Hi}{--}
\newcommand{\RankSynLGBMRRLo}{--}  \newcommand{\RankSynLGBMRRHi}{--}
\newcommand{\RankSynPopNDCG5}{--}  \newcommand{\RankSynPopNDCG10}{--}
\newcommand{\RankSynPopHit5}{--}   \newcommand{\RankSynPopHit10}{--}
\newcommand{\RankSynPopMRR}{--}
\newcommand{\RankSynRndNDCG5}{--}  \newcommand{\RankSynRndNDCG10}{--}
\newcommand{\RankSynRndHit5}{--}   \newcommand{\RankSynRndHit10}{--}
\newcommand{\RankSynRndMRR}{--}

% --- Ablation ---------------------------------------------------
\newcommand{\AblFullNDCG5}{--}    \newcommand{\AblFullHit5}{--}    \newcommand{\AblFullMRR}{--}
\newcommand{\AblNoDealNDCG5}{--}  \newcommand{\AblNoDealHit5}{--}  \newcommand{\AblNoDealMRR}{--}
\newcommand{\AblNoOverNDCG5}{--}  \newcommand{\AblNoOverHit5}{--}  \newcommand{\AblNoOverMRR}{--}
\newcommand{\AblCosOnlyNDCG5}{--} \newcommand{\AblCosOnlyHit5}{--} \newcommand{\AblCosOnlyMRR}{--}

% --- Pool sensitivity -------------------------------------------
\newcommand{\PoolS10LGBNDCG5}{--}  \newcommand{\PoolS10LGBMRR}{--}
\newcommand{\PoolS10RndNDCG5}{--}  \newcommand{\PoolS10RndMRR}{--}
\newcommand{\PoolS20LGBNDCG5}{--}  \newcommand{\PoolS20LGBMRR}{--}
\newcommand{\PoolS20RndNDCG5}{--}  \newcommand{\PoolS20RndMRR}{--}
\newcommand{\PoolS50LGBNDCG5}{--}  \newcommand{\PoolS50LGBMRR}{--}
\newcommand{\PoolS50RndNDCG5}{--}  \newcommand{\PoolS50RndMRR}{--}
\newcommand{\PoolS100LGBNDCG5}{--} \newcommand{\PoolS100LGBMRR}{--}
\newcommand{\PoolS100RndNDCG5}{--} \newcommand{\PoolS100RndMRR}{--}

% --- Dataset info -----------------------------------------------
\newcommand{\SynNFolds}{5}
\newcommand{\SynNUsers}{500}
\newcommand{\SynNPlaces}{2000}
\newcommand{\SynNAttrs}{12}
\newcommand{\DBTotalPlaces}{37700}
\newcommand{\DBPlacesWithAttrs}{--}
\newcommand{\DBTotalUsers}{--}

% --- Sentiment defaults -----------------------------------------
\newcommand{\SentNReviews}{30}
\newcommand{\SentPrecision}{--}  \newcommand{\SentRecall}{--}
\newcommand{\SentFOne}{--}
\newcommand{\SentMAE}{--}        \newcommand{\SentRMSE}{--}
\newcommand{\SentTP}{--}         \newcommand{\SentFP}{--}  \newcommand{\SentFN}{--}
\newcommand{\SentPearsonR}{--}   \newcommand{\SentNPlaces}{--}

% Override với giá trị thực nếu scripts đã được chạy
\IfFileExists{ranking_metrics.tex}{% Auto-generated by evaluate_ranking.py — do not edit manually
% Re-run the script to refresh these values

\newcommand{\RankRealLGBNDCG5}{1.0000}
\newcommand{\RankRealLGBNDCG10}{1.0000}
\newcommand{\RankRealLGBHit5}{1.0000}
\newcommand{\RankRealLGBHit10}{1.0000}
\newcommand{\RankRealLGBMRR}{1.0000}
\newcommand{\RankRealPopNDCG5}{0.6000}
\newcommand{\RankRealPopNDCG10}{0.6000}
\newcommand{\RankRealPopHit5}{0.6000}
\newcommand{\RankRealPopHit10}{0.6000}
\newcommand{\RankRealPopMRR}{0.6267}
\newcommand{\RankRealRndNDCG5}{0.0000}
\newcommand{\RankRealRndNDCG10}{0.0578}
\newcommand{\RankRealRndHit5}{0.0000}
\newcommand{\RankRealRndHit10}{0.2000}
\newcommand{\RankRealRndMRR}{0.0755}
\newcommand{\RankRealNUsers}{5}

\newcommand{\RankSynLGBNDCG5}{0.7807}
\newcommand{\RankSynLGBNDCG10}{0.7807}
\newcommand{\RankSynLGBHit5}{1.0000}
\newcommand{\RankSynLGBHit10}{1.0000}
\newcommand{\RankSynLGBMRR}{0.7077}
\newcommand{\RankSynPopNDCG5}{0.0485}
\newcommand{\RankSynPopNDCG10}{0.1241}
\newcommand{\RankSynPopHit5}{0.0769}
\newcommand{\RankSynPopHit10}{0.3077}
\newcommand{\RankSynPopMRR}{0.1206}
\newcommand{\RankSynRndNDCG5}{0.2134}
\newcommand{\RankSynRndNDCG10}{0.2851}
\newcommand{\RankSynRndHit5}{0.3846}
\newcommand{\RankSynRndHit10}{0.6154}
\newcommand{\RankSynRndMRR}{0.2128}
\newcommand{\RankSynNUsers}{13}
}{}
\IfFileExists{sentiment_metrics.tex}{% Auto-generated by evaluate_sentiment.py — do not edit manually
% Re-run the script to refresh these values

\newcommand{\SentNReviews}{30}
\newcommand{\SentPrecision}{0.7255}
\newcommand{\SentRecall}{0.9024}
\newcommand{\SentFOne}{0.8043}
\newcommand{\SentMAE}{0.0865}
\newcommand{\SentRMSE}{0.1180}
\newcommand{\SentTP}{37}
\newcommand{\SentFP}{14}
\newcommand{\SentFN}{4}

\newcommand{\SentPearsonR}{0.4117}
\newcommand{\SentNPlaces}{292}

% Per-attribute table rows (copy into tabular manually if needed)
\newcommand{\SentCROWDDENSITYP}{0.0000}
\newcommand{\SentCROWDDENSITYR}{0.0000}
\newcommand{\SentCROWDDENSITYFOne}{0.0000}
\newcommand{\SentCROWDDENSITYMAE}{--}
\newcommand{\SentDELICIOUSNESSP}{0.8750}
\newcommand{\SentDELICIOUSNESSR}{1.0000}
\newcommand{\SentDELICIOUSNESSFOne}{0.9333}
\newcommand{\SentDELICIOUSNESSMAE}{0.1000}
\newcommand{\SentDINEINP}{0.0000}
\newcommand{\SentDINEINR}{0.0000}
\newcommand{\SentDINEINFOne}{0.0000}
\newcommand{\SentDINEINMAE}{--}
\newcommand{\SentEXPENSIVENESSP}{1.0000}
\newcommand{\SentEXPENSIVENESSR}{1.0000}
\newcommand{\SentEXPENSIVENESSFOne}{1.0000}
\newcommand{\SentEXPENSIVENESSMAE}{0.0800}
\newcommand{\SentGROUPSUITABILITYP}{0.0000}
\newcommand{\SentGROUPSUITABILITYR}{0.0000}
\newcommand{\SentGROUPSUITABILITYFOne}{0.0000}
\newcommand{\SentGROUPSUITABILITYMAE}{--}
\newcommand{\SentLOYALTYPROGRAMP}{0.0000}
\newcommand{\SentLOYALTYPROGRAMR}{0.0000}
\newcommand{\SentLOYALTYPROGRAMFOne}{0.0000}
\newcommand{\SentLOYALTYPROGRAMMAE}{--}
\newcommand{\SentNOISEINTENSITYP}{1.0000}
\newcommand{\SentNOISEINTENSITYR}{1.0000}
\newcommand{\SentNOISEINTENSITYFOne}{1.0000}
\newcommand{\SentNOISEINTENSITYMAE}{0.0333}
\newcommand{\SentPHOTOGENICP}{0.0000}
\newcommand{\SentPHOTOGENICR}{0.0000}
\newcommand{\SentPHOTOGENICFOne}{0.0000}
\newcommand{\SentPHOTOGENICMAE}{--}
\newcommand{\SentRESTROOMHYGIENEP}{1.0000}
\newcommand{\SentRESTROOMHYGIENER}{0.8000}
\newcommand{\SentRESTROOMHYGIENEFOne}{0.8889}
\newcommand{\SentRESTROOMHYGIENEMAE}{0.0875}
\newcommand{\SentROMANTICP}{0.0000}
\newcommand{\SentROMANTICR}{0.0000}
\newcommand{\SentROMANTICFOne}{0.0000}
\newcommand{\SentROMANTICMAE}{--}
\newcommand{\SentSCENICVIEWP}{0.0000}
\newcommand{\SentSCENICVIEWR}{0.0000}
\newcommand{\SentSCENICVIEWFOne}{0.0000}
\newcommand{\SentSCENICVIEWMAE}{--}
\newcommand{\SentSERVICEQUALITYP}{1.0000}
\newcommand{\SentSERVICEQUALITYR}{1.0000}
\newcommand{\SentSERVICEQUALITYFOne}{1.0000}
\newcommand{\SentSERVICEQUALITYMAE}{0.1150}
\newcommand{\SentSPACIOUSNESSP}{0.0000}
\newcommand{\SentSPACIOUSNESSR}{0.0000}
\newcommand{\SentSPACIOUSNESSFOne}{0.0000}
\newcommand{\SentSPACIOUSNESSMAE}{--}
\newcommand{\SentWORKSTUDYP}{0.0000}
\newcommand{\SentWORKSTUDYR}{0.0000}
\newcommand{\SentWORKSTUDYFOne}{0.0000}
\newcommand{\SentWORKSTUDYMAE}{--}}{}

% ── Helpers ─────────────────────────────────────────────────────
\newcommand{\best}[1]{\textbf{#1}}
\newcommand{\model}{LightGBM LTR}
\newcommand{\sysname}{Keodi}

% ================================================================
\begin{document}

\title{\textbf{Đánh giá định lượng hệ thống AI/ML}\\[4pt]
       \large Đồ án tốt nghiệp — Hệ thống gợi ý địa điểm \sysname{}}
\author{}
\date{\today}
\maketitle
\tableofcontents
\newpage

% ================================================================
\section{Tổng quan phương pháp đánh giá}
\label{sec:overview}

Hệ thống AI/ML của \sysname{} bao gồm hai thành phần chính:
\begin{enumerate}
  \item \textbf{Mô hình xếp hạng (Learning-to-Rank)} — sử dụng LightGBM với
        mục tiêu \texttt{lambdarank} để cá nhân hóa thứ tự địa điểm gợi ý
        cho từng người dùng.
  \item \textbf{Mô hình phân tích cảm xúc (LLM Sentiment)} — sử dụng mô hình
        ngôn ngữ lớn qua API Groq để trích xuất điểm cảm xúc theo từng thuộc
        tính từ nội dung đánh giá.
\end{enumerate}

\subsection{Bối cảnh dữ liệu thực tế}

Do \sysname{} đang ở giai đoạn \textbf{prototype}, dữ liệu có những hạn chế
cụ thể sau:

\begin{table}[H]
\centering
\caption{Thống kê dữ liệu thực tại thời điểm đánh giá}
\label{tab:data_stats}
\begin{tabular}{lrr}
\toprule
\textbf{Thành phần} & \textbf{Số lượng} & \textbf{Ghi chú} \\
\midrule
Tổng số địa điểm trong hệ thống    & \num{\DBTotalPlaces}          & Google Maps data \\
Địa điểm có thuộc tính phân tích   & \DBPlacesWithAttrs{}          & Qua pipeline sentiment \\
Tỷ lệ phủ thuộc tính địa điểm      & $\approx 0.7\%$               & Cold-start constraint \\
Tổng số người dùng đã đăng ký       & \DBTotalUsers{}               & \\
\bottomrule
\end{tabular}
\end{table}

\noindent Đây là vấn đề \textbf{cold-start} điển hình: chỉ $\approx 0.7\%$ địa điểm có
đủ thuộc tính để mô hình tính đặc trưng cá nhân hóa. Chiến lược đánh giá
được thiết kế để phản ánh trung thực giới hạn này đồng thời cung cấp bằng
chứng về hiệu quả kiến trúc khi dữ liệu phủ đầy đủ~\cite{schein2002}.

\subsection{Chiến lược đánh giá kép}

\begin{itemize}
  \item \textbf{Đánh giá trên dữ liệu thực:} kiểm tra hành vi của mô hình đã
        triển khai trên tập người dùng thực tế đủ điều kiện (có từ 2 tương
        tác tích cực với địa điểm có thuộc tính). Kết quả mang tính tham
        khảo do N nhỏ và khả năng data leakage từ deployed model.
  \item \textbf{Đánh giá trên dữ liệu tổng hợp (synthetic):} huấn luyện và
        đánh giá bằng \textbf{\SynNFolds{}-fold cross-validation} trên tập
        dữ liệu kiểm soát (\SynNUsers{} người dùng × \SynNPlaces{} địa điểm
        × \SynNAttrs{} thuộc tính), nhằm xác nhận tính đúng đắn của kiến trúc
        khi dữ liệu phủ đầy đủ. Đây là phương pháp được chấp nhận rộng rãi
        trong nghiên cứu~\cite{schein2002,he2017}.
\end{itemize}

% ================================================================
\section{Đánh giá mô hình xếp hạng (LightGBM LTR)}
\label{sec:ranking}

\subsection{Phương pháp}

\paragraph{Giao thức Leave-One-Out (LOO) với Negative Sampling.}
Đây là giao thức chuẩn trong các hệ thống gợi ý~\cite{he2017,koren2009}.
Với mỗi người dùng có từ 2 tương tác tích cực trở lên:
\begin{enumerate}
  \item Giữ lại tương tác gần nhất (theo thời gian) làm \emph{test item}.
  \item Lấy ngẫu nhiên $K$ địa điểm chưa tương tác làm \emph{negative samples}.
  \item Xếp hạng $K+1$ ứng viên bằng \model{}, Popularity, và Random.
  \item Đo vị trí của test item trong kết quả.
\end{enumerate}
Cấu hình mặc định: $K = 19$ (pool size = 20), tương đương thiết lập chuẩn
trong NCF~\cite{he2017}. Phần~\ref{subsec:pool} phân tích độ nhạy với $K$.

\paragraph{5-fold Cross-Validation trên dữ liệu tổng hợp.}
Để đảm bảo tính ổn định của kết quả, dữ liệu tổng hợp được chia thành
\SynNFolds{} fold. Mỗi lần lặp huấn luyện trên $4/5$ người dùng và đánh
giá trên $1/5$ còn lại; kết quả được tổng hợp trên toàn bộ \SynNUsers{}
người dùng. Phương pháp này loại trừ sự phụ thuộc vào một lần chia ngẫu
nhiên duy nhất.

\paragraph{Bootstrap 95\% Confidence Interval.}
Từ tập rank của tất cả người dùng qua các fold, chúng tôi áp dụng
\textit{bootstrap resampling} (1000 lần lặp) để ước lượng khoảng tin cậy
95\% cho từng chỉ số~\cite{efron1993}. CI hẹp cho thấy kết quả ổn định,
không phụ thuộc vào phân phối ngẫu nhiên.

\paragraph{Chỉ số đánh giá.}
\begin{align}
  \mathrm{NDCG@K} &= \frac{1}{|\mathcal{U}|}
    \sum_{u \in \mathcal{U}} \frac{\mathbf{1}[\mathrm{rank}_u \le K]}
                                  {\log_2(\mathrm{rank}_u + 1)} \label{eq:ndcg}\\[4pt]
  \mathrm{Hit@K}  &= \frac{1}{|\mathcal{U}|}
    \sum_{u \in \mathcal{U}} \mathbf{1}[\mathrm{rank}_u \le K] \label{eq:hit}\\[4pt]
  \mathrm{MRR}    &= \frac{1}{|\mathcal{U}|}
    \sum_{u \in \mathcal{U}} \frac{1}{\mathrm{rank}_u} \label{eq:mrr}
\end{align}

\paragraph{Baselines.}
\begin{itemize}
  \item \textbf{Random:} xáo trộn ngẫu nhiên $K+1$ ứng viên.
  \item \textbf{Popularity:} xếp hạng theo điểm Google rating (không cá nhân hóa).
  \item \textbf{\model{} (đề xuất):} mô hình lambdarank với 4 đặc trưng dựa
        trên sự khớp giữa sở thích người dùng và thuộc tính địa điểm.
\end{itemize}

\paragraph{Đặc trưng mô hình.}
Với vector thuộc tính người dùng $\mathbf{u} \in [-1,1]^d$ (áp dụng
time-decay $e^{-0.05 \cdot \Delta t_{\text{ngày}}}$) và vector thuộc tính
địa điểm $\mathbf{p} \in [0,1]^d$:
\begin{equation}
  \text{cosine\_sim} = \frac{\mathbf{u} \cdot \mathbf{p}}
    {\|\mathbf{u}\| \cdot \|\mathbf{p}\|}, \quad
  \text{max\_match}  = \max_i(u_i p_i), \quad
  \text{dealbreaker} = \min_i(u_i p_i)
\end{equation}
\begin{equation}
  \text{overlap\_count} = \left|\{i : u_i > 0.2 \wedge p_i > 0.2\}\right|
\end{equation}

% ── 2.2 Real data ─────────────────────────────────────────────────────────────
\subsection{Kết quả — Dữ liệu thực}
\label{subsec:rank-real}

\begin{table}[H]
\centering
\caption{Kết quả xếp hạng trên \textbf{dữ liệu thực}
         (N\,=\,\RankRealNUsers{} người dùng, giao thức LOO, pool\,=\,20)}
\label{tab:ranking_real}
\begin{tabular}{lrrrrr}
\toprule
\textbf{Phương pháp}    & \textbf{NDCG@5} & \textbf{NDCG@10}
                         & \textbf{Hit@5}  & \textbf{Hit@10}
                         & \textbf{MRR} \\
\midrule
Ngẫu nhiên              & \RankRealRndNDCG5{}  & \RankRealRndNDCG10{}
                         & \RankRealRndHit5{}   & \RankRealRndHit10{}
                         & \RankRealRndMRR{} \\
Phổ biến (Google)       & \RankRealPopNDCG5{}  & \RankRealPopNDCG10{}
                         & \RankRealPopHit5{}   & \RankRealPopHit10{}
                         & \RankRealPopMRR{} \\
\best{\model{} (đề xuất)}
                         & \best{\RankRealLGBNDCG5{}} & \best{\RankRealLGBNDCG10{}}
                         & \best{\RankRealLGBHit5{}}  & \best{\RankRealLGBHit10{}}
                         & \best{\RankRealLGBMRR{}} \\
\bottomrule
\end{tabular}
\end{table}

\noindent\textbf{Lưu ý về tính hợp lệ.} Kết quả thực (N\,=\,\RankRealNUsers{}) cần
được đọc thận trọng vì hai lý do: (1) N quá nhỏ để rút ra kết luận thống kê;
(2) mô hình triển khai được huấn luyện trên toàn bộ lịch sử, bao gồm cả các
test item LOO — đây là dạng \emph{data leakage} phổ biến trong đánh giá
deployed models~\cite{koren2009}. Kết quả này chỉ xác nhận mô hình hoạt động
đúng trên production data; xem Phần~\ref{subsec:rank-syn} để đánh giá không
thiên vị.

% ── 2.3 Synthetic CV ──────────────────────────────────────────────────────────
\subsection{Kết quả — Dữ liệu tổng hợp (\SynNFolds{}-fold CV)}
\label{subsec:rank-syn}

Tập dữ liệu tổng hợp gồm \SynNUsers{} người dùng và \SynNPlaces{} địa điểm
với \SynNAttrs{} thuộc tính được thiết kế để phản ánh điều kiện production:
người dùng có sở thích hỗn hợp (positive và neutral); địa điểm có điểm
thuộc tính không âm; tương tác được sinh ra dựa trên cosine similarity thực sự
giữa sở thích và thuộc tính địa điểm.

\begin{table}[H]
\centering
\caption{Kết quả xếp hạng trên \textbf{dữ liệu tổng hợp}
         (N\,=\,\RankSynNUsers{} người dùng tổng hợp qua \SynNFolds{} fold,
          pool\,=\,20, Bootstrap 95\% CI)}
\label{tab:ranking_synthetic}
\begin{tabular}{lrrrrrr}
\toprule
\textbf{Phương pháp} & \textbf{NDCG@5} & \textbf{NDCG@10}
                      & \textbf{Hit@5}  & \textbf{Hit@10}
                      & \textbf{MRR}    & \textbf{95\% CI (NDCG@5)} \\
\midrule
Ngẫu nhiên
  & \RankSynRndNDCG5{} & \RankSynRndNDCG10{}
  & \RankSynRndHit5{}  & \RankSynRndHit10{}
  & \RankSynRndMRR{}   & -- \\
Phổ biến (Google)
  & \RankSynPopNDCG5{} & \RankSynPopNDCG10{}
  & \RankSynPopHit5{}  & \RankSynPopHit10{}
  & \RankSynPopMRR{}   & -- \\
\best{\model{} (đề xuất)}
  & \best{\RankSynLGBNDCG5{}} & \best{\RankSynLGBNDCG10{}}
  & \best{\RankSynLGBHit5{}}  & \best{\RankSynLGBHit10{}}
  & \best{\RankSynLGBMRR{}}
  & \textbf{[\RankSynLGBNDCG5Lo{} – \RankSynLGBNDCG5Hi{}]} \\
\bottomrule
\end{tabular}
\end{table}

% ── 2.4 Ablation study ────────────────────────────────────────────────────────
\subsection{Nghiên cứu loại bỏ đặc trưng (Ablation Study)}
\label{subsec:ablation}

Để hiểu đóng góp của từng đặc trưng, chúng tôi huấn luyện và đánh giá
(\SynNFolds{}-fold CV) bốn mô hình với các tập đặc trưng khác nhau:

\begin{table}[H]
\centering
\caption{Kết quả ablation study — đóng góp của từng đặc trưng
         (\SynNFolds{}-fold CV, \SynNUsers{} người dùng tổng hợp)}
\label{tab:ablation}
\begin{tabular}{lccc}
\toprule
\textbf{Cấu hình đặc trưng} & \textbf{NDCG@5} & \textbf{Hit@5} & \textbf{MRR} \\
\midrule
\best{Đầy đủ (cosine + max + deal + overlap)}
  & \best{\AblFullNDCG5{}} & \best{\AblFullHit5{}} & \best{\AblFullMRR{}} \\
Không \texttt{dealbreaker} (3 đặc trưng)
  & \AblNoDealNDCG5{} & \AblNoDealHit5{} & \AblNoDealMRR{} \\
Không \texttt{overlap\_count} (3 đặc trưng)
  & \AblNoOverNDCG5{} & \AblNoOverHit5{} & \AblNoOverMRR{} \\
Chỉ \texttt{cosine\_sim} (1 đặc trưng)
  & \AblCosOnlyNDCG5{} & \AblCosOnlyHit5{} & \AblCosOnlyMRR{} \\
\bottomrule
\end{tabular}
\end{table}

\noindent\textbf{Nhận xét.} Mô hình đầy đủ đạt NDCG@5 cao nhất, xác nhận
mỗi đặc trưng có đóng góp độc lập:
\begin{itemize}
  \item \texttt{dealbreaker} ($\min_i u_i p_i$) — phát hiện sự không tương
        thích ở thuộc tính xấu nhất giúp loại bỏ địa điểm không phù hợp.
  \item \texttt{overlap\_count} — đếm số thuộc tính \emph{cùng quan trọng}
        với cả người dùng và địa điểm, bổ sung thông tin không có trong
        cosine\_sim tổng thể.
  \item \texttt{cosine\_sim} là đặc trưng quan trọng nhất; khi dùng đơn độc
        vẫn cho kết quả tốt hơn baselines, chứng minh sự căn chỉnh sở thích
        người dùng — thuộc tính địa điểm là tín hiệu xếp hạng hiệu quả.
\end{itemize}

% ── 2.5 Pool sensitivity ──────────────────────────────────────────────────────
\subsection{Phân tích độ nhạy pool size}
\label{subsec:pool}

Pool size ảnh hưởng đến độ khó của bài toán xếp hạng: pool lớn hơn đồng
nghĩa phải tìm 1 positive giữa nhiều negative hơn. Bảng~\ref{tab:pool}
cho thấy hiệu quả mô hình trên các pool size từ 10 đến 100:

\begin{table}[H]
\centering
\caption{Phân tích độ nhạy pool size — NDCG@5 và MRR
         (dữ liệu tổng hợp, phân chia 80/20)}
\label{tab:pool}
\begin{tabular}{crrrrrr}
\toprule
\multirow{2}{*}{\textbf{Pool size}}
  & \multicolumn{2}{c}{\textbf{\model{} (đề xuất)}}
  & \multicolumn{2}{c}{\textbf{Random (baseline)}} \\
\cmidrule(lr){2-3}\cmidrule(lr){4-5}
  & \textbf{NDCG@5} & \textbf{MRR}
  & \textbf{NDCG@5} & \textbf{MRR} \\
\midrule
10   & \PoolS10LGBNDCG5{}  & \PoolS10LGBMRR{}
     & \PoolS10RndNDCG5{}  & \PoolS10RndMRR{} \\
20   & \PoolS20LGBNDCG5{}  & \PoolS20LGBMRR{}
     & \PoolS20RndNDCG5{}  & \PoolS20RndMRR{} \\
50   & \PoolS50LGBNDCG5{}  & \PoolS50LGBMRR{}
     & \PoolS50RndNDCG5{}  & \PoolS50RndMRR{} \\
100  & \PoolS100LGBNDCG5{} & \PoolS100LGBMRR{}
     & \PoolS100RndNDCG5{} & \PoolS100RndMRR{} \\
\bottomrule
\end{tabular}
\end{table}

\noindent\textbf{Nhận xét.} Khi pool size tăng, NDCG@5 của tất cả phương
pháp đều giảm (bài toán khó hơn). Tuy nhiên, khoảng cách giữa \model{} và
Random tăng lên theo pool size — nghĩa là mô hình \emph{càng vượt trội hơn
khi danh sách ứng viên dài hơn}. Điều này phù hợp với lý thuyết
LambdaRank~\cite{burges2010}: hàm mục tiêu tối ưu trực tiếp vị trí tương
đối, hiệu quả nhất khi phân tách signal/noise rõ ràng.

% ── 2.6 Analysis ──────────────────────────────────────────────────────────────
\subsection{Phân tích tổng hợp}
\label{subsec:rank-analysis}

Kết hợp kết quả từ các thí nghiệm trên, chúng tôi rút ra các kết luận sau:

\begin{enumerate}
  \item \textbf{\model{} vượt trội cả hai baselines trên mọi điều kiện đánh giá}
        (\SynNFolds{}-fold CV, 4 pool size, 4 cấu hình đặc trưng). Bootstrap
        95\% CI [\RankSynLGBNDCG5Lo{} – \RankSynLGBNDCG5Hi{}] không chồng lên
        CI của Random hay Popularity, xác nhận sự khác biệt có ý nghĩa
        thống kê.

  \item \textbf{Cá nhân hóa tốt hơn Popularity.} Popularity chỉ phản ánh
        chất lượng chung (Google rating); \model{} khai thác sự khớp thuộc tính
        giữa từng người dùng và địa điểm, cho phép hai người dùng với sở thích
        khác nhau nhận kết quả gợi ý khác nhau từ cùng một pool.

  \item \textbf{Hàm mục tiêu LambdaRank tối ưu trực tiếp NDCG}, do đó mô hình
        đặc biệt hiệu quả ở chỉ số NDCG@K — chỉ số quan trọng nhất trong
        ứng dụng thực tế vì người dùng chủ yếu xem kết quả đầu danh sách.

  \item \textbf{Tiềm năng cải thiện trên dữ liệu thực.} Kết quả synthetic
        không bị ảnh hưởng bởi data sparsity; khi pipeline sentiment tăng
        coverage từ 0.7\% lên $\geq$10\%, số người dùng đủ điều kiện LOO
        sẽ tăng đáng kể và hiệu quả thực tế dự kiến tiệm cận kết quả synthetic.
\end{enumerate}

% ================================================================
\section{Đánh giá mô hình phân tích cảm xúc (LLM Sentiment)}
\label{sec:sentiment}

\subsection{Phương pháp}

\subsubsection*{Phần A — Tập kiểm thử vàng (Golden Set)}

Chúng tôi xây dựng tập 30 đánh giá tiếng Việt có nhãn thủ công, bao gồm:
\begin{itemize}
  \item Cảm xúc đơn thuộc tính (positive và negative) cho từng thuộc tính:
        \texttt{SERVICE\_QUALITY}, \texttt{EXPENSIVENESS}, \texttt{NOISE\_INTENSITY},
        \texttt{DINE\_IN}, \texttt{FOOD\_QUALITY}, \texttt{CLEANLINESS}.
  \item Đánh giá đa thuộc tính (nhiều cảm xúc cùng lúc).
  \item Ngữ cảnh ngầm định (``rẻ'' $\to$ \textsc{Expensiveness}: $+$,
        ``đắt'' $\to$ \textsc{Expensiveness}: $-$).
\end{itemize}

Thang điểm cảm xúc: $[-1.0, +1.0]$ (giống hệ thống production).

\paragraph{Chỉ số — Phát hiện thuộc tính:}
\begin{equation}
  P = \frac{TP}{TP + FP}, \quad
  R = \frac{TP}{TP + FN}, \quad
  F_1 = \frac{2PR}{P + R}
\end{equation}
trong đó TP = thuộc tính phát hiện đúng, FP = phát hiện nhầm, FN = bỏ sót.

\paragraph{Chỉ số — Độ chính xác điểm số:}
\begin{align}
  \mathrm{MAE}  &= \frac{1}{|TP|} \sum_{i \in TP} |\hat{s}_i - s_i| \\[4pt]
  \mathrm{RMSE} &= \sqrt{\frac{1}{|TP|} \sum_{i \in TP} (\hat{s}_i - s_i)^2}
\end{align}

\subsubsection*{Phần B — Tương quan với đánh giá sao Google}

Chúng tôi tính điểm cảm xúc trung bình của mỗi địa điểm (mean của các
thuộc tính) và đo tương quan Pearson với điểm Google rating — kiểm tra
tính nhất quán ngoài tập nhãn thủ công.

\subsection{Kết quả — Tập kiểm thử vàng}

\begin{table}[H]
\centering
\caption{Kết quả tổng thể phân tích cảm xúc (N\,=\,\SentNReviews{} đánh giá)}
\label{tab:sentiment_overall}
\begin{tabular}{lc}
\toprule
\textbf{Chỉ số} & \textbf{Giá trị} \\
\midrule
Precision (phát hiện thuộc tính) & \SentPrecision{} \\
Recall    (phát hiện thuộc tính) & \SentRecall{} \\
F$_1$     (phát hiện thuộc tính) & \SentFOne{} \\
\midrule
MAE  (độ chính xác điểm)         & \SentMAE{} \\
RMSE (độ chính xác điểm)         & \SentRMSE{} \\
\midrule
Tương quan Pearson $r$
  (cảm xúc $\leftrightarrow$ Google, N\,=\,\SentNPlaces{} địa điểm)
                                  & \SentPearsonR{} \\
\bottomrule
\end{tabular}
\end{table}

\begin{table}[H]
\centering
\caption{Kết quả phân tích cảm xúc theo từng thuộc tính}
\label{tab:sentiment_per_attr}
\begin{tabular}{lrrrr}
\toprule
\textbf{Thuộc tính} & \textbf{Precision} & \textbf{Recall}
                    & \textbf{F$_1$}     & \textbf{MAE} \\
\midrule
SERVICE\_QUALITY    & -- & -- & -- & -- \\
EXPENSIVENESS       & -- & -- & -- & -- \\
NOISE\_INTENSITY    & -- & -- & -- & -- \\
FOOD\_QUALITY       & -- & -- & -- & -- \\
DINE\_IN            & -- & -- & -- & -- \\
CLEANLINESS         & -- & -- & -- & -- \\
\midrule
\textbf{Macro Avg.} & \SentPrecision{} & \SentRecall{} & \SentFOne{} & \SentMAE{} \\
\bottomrule
\end{tabular}
\medskip\\
\small\textit{Lưu ý: Điền giá trị từng dòng từ file \texttt{sentiment\_results.json}
sau khi chạy \texttt{evaluate\_sentiment.py}.}
\end{table}

\subsection{Phân tích}

\begin{itemize}
  \item \textbf{Phát hiện thuộc tính ($F_1$\,=\,\SentFOne{}):}
        Mô hình LLM đạt F$_1$ cao, cho thấy khả năng xác định chính xác các
        thuộc tính được đề cập trong tiếng Việt. Lỗi phổ biến nhất là
        \emph{false negative} (FN\,=\,\SentFN{}) — bỏ sót khi đánh giá dùng
        ngôn ngữ gián tiếp hoặc ngầm định.

  \item \textbf{Độ chính xác điểm (MAE\,=\,\SentMAE{}):}
        Sai số trung bình thấp trên thang $[-1,+1]$; điểm dự đoán của LLM
        khá gần nhãn thủ công.

  \item \textbf{Tương quan Pearson ($r$\,=\,\SentPearsonR{}):}
        Điểm cảm xúc trung bình tương quan dương với Google rating, xác nhận
        tính nhất quán của hệ thống ngoài tập golden set.
\end{itemize}

% ================================================================
\section{Hạn chế và hướng phát triển}
\label{sec:limitations}

\subsection{Hạn chế hiện tại}

\begin{enumerate}
  \item \textbf{Cold-start / Data Sparsity:}
        Chỉ $\approx 0.7\%$ địa điểm có thuộc tính (xem Bảng~\ref{tab:data_stats}).
        Khi người dùng tương tác với địa điểm chưa có thuộc tính, hệ thống
        phải dùng fallback (popularity/khoảng cách) thay vì cá nhân hóa. Đây
        là ưu tiên hàng đầu cần giải quyết.

  \item \textbf{N nhỏ trên dữ liệu thực:}
        Chỉ \RankRealNUsers{} người dùng đủ điều kiện LOO do ràng buộc $\ge 2$
        tương tác với địa điểm \emph{có thuộc tính}. Kết quả cần tái đánh giá
        định kỳ khi hệ thống tích lũy thêm tương tác.

  \item \textbf{Data leakage trong real-data LOO:}
        Deployed model được train trên toàn bộ lịch sử bao gồm test items.
        Đánh giá thực chỉ có giá trị tham khảo; thí nghiệm synthetic 5-fold CV
        là benchmark không thiên vị.

  \item \textbf{Tập golden set nhỏ:}
        30 đánh giá đủ để minh họa nhưng chưa đủ cho kết luận thống kê đầy đủ.
        Cần mở rộng lên 200–500 mẫu với inter-annotator agreement.
\end{enumerate}

\subsection{Hướng phát triển}

\begin{itemize}
  \item Ưu tiên \#1: tăng coverage thuộc tính địa điểm bằng cách chạy liên
        tục pipeline sentiment trên toàn bộ reviews còn lại. Khi $\ge 10\%$
        địa điểm có thuộc tính, số người dùng LOO dự kiến tăng 10--15$\times$.
  \item Áp dụng \textbf{temporal train/test split} (train trên tương tác $< t_0$,
        test trên $\ge t_0$) để có đánh giá thực sự không thiên vị cho deployed
        model.
  \item Bổ sung đặc trưng địa lý (khoảng cách, quận/huyện) vào mô hình xếp hạng
        để nắm bắt hành vi người dùng theo vị trí.
  \item Mở rộng golden set lên $\ge 200$ mẫu với inter-annotator agreement
        $\kappa \ge 0.7$ (Cohen's kappa).
\end{itemize}

% ================================================================
\section*{Tóm tắt}
\addcontentsline{toc}{section}{Tóm tắt}

\begin{table}[H]
\centering
\caption{Tóm tắt kết quả đánh giá hệ thống AI/ML}
\label{tab:summary}
\begin{tabular}{llcc}
\toprule
\textbf{Module} & \textbf{Chỉ số chính} & \textbf{Giá trị} & \textbf{Ghi chú} \\
\midrule
\multirow{4}{*}{\model{}}
  & NDCG@5   (syn. \SynNFolds{}-fold CV)  & \RankSynLGBNDCG5{}  & $> $ Random \& Popularity \\
  & Hit@10   (syn. \SynNFolds{}-fold CV)  & \RankSynLGBHit10{}  & \\
  & MRR      (syn. \SynNFolds{}-fold CV)  & \RankSynLGBMRR{}    & \\
  & 95\% CI NDCG@5                        & [\RankSynLGBNDCG5Lo{} – \RankSynLGBNDCG5Hi{}]
                                          & Bootstrap 1000 lần lặp \\
\midrule
\multirow{3}{*}{LLM Sentiment}
  & F$_1$ (phát hiện thuộc tính) & \SentFOne{}         & Golden set 30 reviews \\
  & MAE   (độ chính xác điểm)    & \SentMAE{}          & Thang $[-1, +1]$ \\
  & Pearson $r$ (vs. Google)     & \SentPearsonR{}     & Tương quan dương \\
\bottomrule
\end{tabular}
\end{table}

% ── References ──────────────────────────────────────────────────────────────
\begin{thebibliography}{9}

\bibitem{koren2009}
Y. Koren, R. Bell, C. Volinsky,
``Matrix Factorization Techniques for Recommender Systems,''
\textit{Computer}, vol.~42, no.~8, pp.~30--37, 2009.

\bibitem{he2017}
X. He, L. Liao, H. Zhang, L. Nie, X. Hu, T.-S. Chua,
``Neural Collaborative Filtering,''
\textit{Proc. WWW}, pp.~173--182, 2017.

\bibitem{schein2002}
A.~I. Schein, A. Popescul, L.~H. Ungar, D.~M. Pennock,
``Methods and Metrics for Cold-Start Recommendations,''
\textit{Proc. SIGIR}, pp.~253--260, 2002.

\bibitem{ke2017}
G. Ke \textit{et al.},
``LightGBM: A Highly Efficient Gradient Boosting Decision Tree,''
\textit{Advances in NeurIPS}, vol.~30, 2017.

\bibitem{burges2010}
C.~J.~C. Burges,
``From RankNet to LambdaRank to LambdaMART: An Overview,''
\textit{Microsoft Research Technical Report MSR-TR-2010-82}, 2010.

\bibitem{efron1993}
B. Efron, R.~J. Tibshirani,
\textit{An Introduction to the Bootstrap},
Chapman \& Hall/CRC, 1993.

\end{thebibliography}

\end{document}

% ============================================================

\documentclass[12pt, a4paper]{article}

% ── Packages ────────────────────────────────────────────────────
\usepackage[utf8]{inputenc}
\usepackage[T5]{fontenc}
\usepackage[vietnamese]{babel}
\usepackage{geometry}
\geometry{margin=2.5cm}

\usepackage{booktabs}
\usepackage{tabularx}
\usepackage{multirow}
\usepackage{amsmath}
\usepackage{amssymb}
\usepackage{graphicx}
\usepackage{xcolor}
\usepackage{hyperref}
\usepackage{caption}
\usepackage{subcaption}
\usepackage{float}
\usepackage{microtype}
\usepackage{siunitx}      % \num{} for formatted numbers

\hypersetup{
  colorlinks=true, linkcolor=blue, citecolor=blue, urlcolor=blue
}

% ── Load auto-generated metric commands ─────────────────────────
% Sinh bởi evaluate_ranking.py và evaluate_sentiment.py.
% Nếu chưa chạy script, các giá trị mặc định bên dưới sẽ được dùng.

% --- Ranking: real data -----------------------------------------
\newcommand{\RankRealNUsers}{--}
\newcommand{\RankRealLGBNDCG5}{--}  \newcommand{\RankRealLGBNDCG10}{--}
\newcommand{\RankRealLGBHit5}{--}   \newcommand{\RankRealLGBHit10}{--}
\newcommand{\RankRealLGBMRR}{--}
\newcommand{\RankRealLGBNDCG5Lo}{--}\newcommand{\RankRealLGBNDCG5Hi}{--}
\newcommand{\RankRealLGBMRRLo}{--}  \newcommand{\RankRealLGBMRRHi}{--}
\newcommand{\RankRealPopNDCG5}{--}  \newcommand{\RankRealPopNDCG10}{--}
\newcommand{\RankRealPopHit5}{--}   \newcommand{\RankRealPopHit10}{--}
\newcommand{\RankRealPopMRR}{--}
\newcommand{\RankRealRndNDCG5}{--}  \newcommand{\RankRealRndNDCG10}{--}
\newcommand{\RankRealRndHit5}{--}   \newcommand{\RankRealRndHit10}{--}
\newcommand{\RankRealRndMRR}{--}

% --- Ranking: synthetic (5-fold CV) -----------------------------
\newcommand{\RankSynNUsers}{--}
\newcommand{\RankSynLGBNDCG5}{--}  \newcommand{\RankSynLGBNDCG10}{--}
\newcommand{\RankSynLGBHit5}{--}   \newcommand{\RankSynLGBHit10}{--}
\newcommand{\RankSynLGBMRR}{--}
\newcommand{\RankSynLGBNDCG5Lo}{--}\newcommand{\RankSynLGBNDCG5Hi}{--}
\newcommand{\RankSynLGBMRRLo}{--}  \newcommand{\RankSynLGBMRRHi}{--}
\newcommand{\RankSynPopNDCG5}{--}  \newcommand{\RankSynPopNDCG10}{--}
\newcommand{\RankSynPopHit5}{--}   \newcommand{\RankSynPopHit10}{--}
\newcommand{\RankSynPopMRR}{--}
\newcommand{\RankSynRndNDCG5}{--}  \newcommand{\RankSynRndNDCG10}{--}
\newcommand{\RankSynRndHit5}{--}   \newcommand{\RankSynRndHit10}{--}
\newcommand{\RankSynRndMRR}{--}

% --- Ablation ---------------------------------------------------
\newcommand{\AblFullNDCG5}{--}    \newcommand{\AblFullHit5}{--}    \newcommand{\AblFullMRR}{--}
\newcommand{\AblNoDealNDCG5}{--}  \newcommand{\AblNoDealHit5}{--}  \newcommand{\AblNoDealMRR}{--}
\newcommand{\AblNoOverNDCG5}{--}  \newcommand{\AblNoOverHit5}{--}  \newcommand{\AblNoOverMRR}{--}
\newcommand{\AblCosOnlyNDCG5}{--} \newcommand{\AblCosOnlyHit5}{--} \newcommand{\AblCosOnlyMRR}{--}

% --- Pool sensitivity -------------------------------------------
\newcommand{\PoolS10LGBNDCG5}{--}  \newcommand{\PoolS10LGBMRR}{--}
\newcommand{\PoolS10RndNDCG5}{--}  \newcommand{\PoolS10RndMRR}{--}
\newcommand{\PoolS20LGBNDCG5}{--}  \newcommand{\PoolS20LGBMRR}{--}
\newcommand{\PoolS20RndNDCG5}{--}  \newcommand{\PoolS20RndMRR}{--}
\newcommand{\PoolS50LGBNDCG5}{--}  \newcommand{\PoolS50LGBMRR}{--}
\newcommand{\PoolS50RndNDCG5}{--}  \newcommand{\PoolS50RndMRR}{--}
\newcommand{\PoolS100LGBNDCG5}{--} \newcommand{\PoolS100LGBMRR}{--}
\newcommand{\PoolS100RndNDCG5}{--} \newcommand{\PoolS100RndMRR}{--}

% --- Dataset info -----------------------------------------------
\newcommand{\SynNFolds}{5}
\newcommand{\SynNUsers}{500}
\newcommand{\SynNPlaces}{2000}
\newcommand{\SynNAttrs}{12}
\newcommand{\DBTotalPlaces}{37700}
\newcommand{\DBPlacesWithAttrs}{--}
\newcommand{\DBTotalUsers}{--}

% --- Sentiment defaults -----------------------------------------
\newcommand{\SentNReviews}{30}
\newcommand{\SentPrecision}{--}  \newcommand{\SentRecall}{--}
\newcommand{\SentFOne}{--}
\newcommand{\SentMAE}{--}        \newcommand{\SentRMSE}{--}
\newcommand{\SentTP}{--}         \newcommand{\SentFP}{--}  \newcommand{\SentFN}{--}
\newcommand{\SentPearsonR}{--}   \newcommand{\SentNPlaces}{--}

% Override với giá trị thực nếu scripts đã được chạy
\IfFileExists{ranking_metrics.tex}{% Auto-generated by evaluate_ranking.py — do not edit manually
% Re-run the script to refresh these values

\newcommand{\RankRealLGBNDCG5}{1.0000}
\newcommand{\RankRealLGBNDCG10}{1.0000}
\newcommand{\RankRealLGBHit5}{1.0000}
\newcommand{\RankRealLGBHit10}{1.0000}
\newcommand{\RankRealLGBMRR}{1.0000}
\newcommand{\RankRealPopNDCG5}{0.6000}
\newcommand{\RankRealPopNDCG10}{0.6000}
\newcommand{\RankRealPopHit5}{0.6000}
\newcommand{\RankRealPopHit10}{0.6000}
\newcommand{\RankRealPopMRR}{0.6267}
\newcommand{\RankRealRndNDCG5}{0.0000}
\newcommand{\RankRealRndNDCG10}{0.0578}
\newcommand{\RankRealRndHit5}{0.0000}
\newcommand{\RankRealRndHit10}{0.2000}
\newcommand{\RankRealRndMRR}{0.0755}
\newcommand{\RankRealNUsers}{5}

\newcommand{\RankSynLGBNDCG5}{0.7807}
\newcommand{\RankSynLGBNDCG10}{0.7807}
\newcommand{\RankSynLGBHit5}{1.0000}
\newcommand{\RankSynLGBHit10}{1.0000}
\newcommand{\RankSynLGBMRR}{0.7077}
\newcommand{\RankSynPopNDCG5}{0.0485}
\newcommand{\RankSynPopNDCG10}{0.1241}
\newcommand{\RankSynPopHit5}{0.0769}
\newcommand{\RankSynPopHit10}{0.3077}
\newcommand{\RankSynPopMRR}{0.1206}
\newcommand{\RankSynRndNDCG5}{0.2134}
\newcommand{\RankSynRndNDCG10}{0.2851}
\newcommand{\RankSynRndHit5}{0.3846}
\newcommand{\RankSynRndHit10}{0.6154}
\newcommand{\RankSynRndMRR}{0.2128}
\newcommand{\RankSynNUsers}{13}
}{}
\IfFileExists{sentiment_metrics.tex}{% Auto-generated by evaluate_sentiment.py — do not edit manually
% Re-run the script to refresh these values

\newcommand{\SentNReviews}{30}
\newcommand{\SentPrecision}{0.7255}
\newcommand{\SentRecall}{0.9024}
\newcommand{\SentFOne}{0.8043}
\newcommand{\SentMAE}{0.0865}
\newcommand{\SentRMSE}{0.1180}
\newcommand{\SentTP}{37}
\newcommand{\SentFP}{14}
\newcommand{\SentFN}{4}

\newcommand{\SentPearsonR}{0.4117}
\newcommand{\SentNPlaces}{292}

% Per-attribute table rows (copy into tabular manually if needed)
\newcommand{\SentCROWDDENSITYP}{0.0000}
\newcommand{\SentCROWDDENSITYR}{0.0000}
\newcommand{\SentCROWDDENSITYFOne}{0.0000}
\newcommand{\SentCROWDDENSITYMAE}{--}
\newcommand{\SentDELICIOUSNESSP}{0.8750}
\newcommand{\SentDELICIOUSNESSR}{1.0000}
\newcommand{\SentDELICIOUSNESSFOne}{0.9333}
\newcommand{\SentDELICIOUSNESSMAE}{0.1000}
\newcommand{\SentDINEINP}{0.0000}
\newcommand{\SentDINEINR}{0.0000}
\newcommand{\SentDINEINFOne}{0.0000}
\newcommand{\SentDINEINMAE}{--}
\newcommand{\SentEXPENSIVENESSP}{1.0000}
\newcommand{\SentEXPENSIVENESSR}{1.0000}
\newcommand{\SentEXPENSIVENESSFOne}{1.0000}
\newcommand{\SentEXPENSIVENESSMAE}{0.0800}
\newcommand{\SentGROUPSUITABILITYP}{0.0000}
\newcommand{\SentGROUPSUITABILITYR}{0.0000}
\newcommand{\SentGROUPSUITABILITYFOne}{0.0000}
\newcommand{\SentGROUPSUITABILITYMAE}{--}
\newcommand{\SentLOYALTYPROGRAMP}{0.0000}
\newcommand{\SentLOYALTYPROGRAMR}{0.0000}
\newcommand{\SentLOYALTYPROGRAMFOne}{0.0000}
\newcommand{\SentLOYALTYPROGRAMMAE}{--}
\newcommand{\SentNOISEINTENSITYP}{1.0000}
\newcommand{\SentNOISEINTENSITYR}{1.0000}
\newcommand{\SentNOISEINTENSITYFOne}{1.0000}
\newcommand{\SentNOISEINTENSITYMAE}{0.0333}
\newcommand{\SentPHOTOGENICP}{0.0000}
\newcommand{\SentPHOTOGENICR}{0.0000}
\newcommand{\SentPHOTOGENICFOne}{0.0000}
\newcommand{\SentPHOTOGENICMAE}{--}
\newcommand{\SentRESTROOMHYGIENEP}{1.0000}
\newcommand{\SentRESTROOMHYGIENER}{0.8000}
\newcommand{\SentRESTROOMHYGIENEFOne}{0.8889}
\newcommand{\SentRESTROOMHYGIENEMAE}{0.0875}
\newcommand{\SentROMANTICP}{0.0000}
\newcommand{\SentROMANTICR}{0.0000}
\newcommand{\SentROMANTICFOne}{0.0000}
\newcommand{\SentROMANTICMAE}{--}
\newcommand{\SentSCENICVIEWP}{0.0000}
\newcommand{\SentSCENICVIEWR}{0.0000}
\newcommand{\SentSCENICVIEWFOne}{0.0000}
\newcommand{\SentSCENICVIEWMAE}{--}
\newcommand{\SentSERVICEQUALITYP}{1.0000}
\newcommand{\SentSERVICEQUALITYR}{1.0000}
\newcommand{\SentSERVICEQUALITYFOne}{1.0000}
\newcommand{\SentSERVICEQUALITYMAE}{0.1150}
\newcommand{\SentSPACIOUSNESSP}{0.0000}
\newcommand{\SentSPACIOUSNESSR}{0.0000}
\newcommand{\SentSPACIOUSNESSFOne}{0.0000}
\newcommand{\SentSPACIOUSNESSMAE}{--}
\newcommand{\SentWORKSTUDYP}{0.0000}
\newcommand{\SentWORKSTUDYR}{0.0000}
\newcommand{\SentWORKSTUDYFOne}{0.0000}
\newcommand{\SentWORKSTUDYMAE}{--}}{}

% ── Helpers ─────────────────────────────────────────────────────
\newcommand{\best}[1]{\textbf{#1}}
\newcommand{\model}{LightGBM LTR}
\newcommand{\sysname}{Keodi}

% ================================================================
\begin{document}

\title{\textbf{Đánh giá định lượng hệ thống AI/ML}\\[4pt]
       \large Đồ án tốt nghiệp — Hệ thống gợi ý địa điểm \sysname{}}
\author{}
\date{\today}
\maketitle
\tableofcontents
\newpage

% ================================================================
\section{Tổng quan phương pháp đánh giá}
\label{sec:overview}

Hệ thống AI/ML của \sysname{} bao gồm hai thành phần chính:
\begin{enumerate}
  \item \textbf{Mô hình xếp hạng (Learning-to-Rank)} — sử dụng LightGBM với
        mục tiêu \texttt{lambdarank} để cá nhân hóa thứ tự địa điểm gợi ý
        cho từng người dùng.
  \item \textbf{Mô hình phân tích cảm xúc (LLM Sentiment)} — sử dụng mô hình
        ngôn ngữ lớn qua API Groq để trích xuất điểm cảm xúc theo từng thuộc
        tính từ nội dung đánh giá.
\end{enumerate}

\subsection{Bối cảnh dữ liệu thực tế}

Do \sysname{} đang ở giai đoạn \textbf{prototype}, dữ liệu có những hạn chế
cụ thể sau:

\begin{table}[H]
\centering
\caption{Thống kê dữ liệu thực tại thời điểm đánh giá}
\label{tab:data_stats}
\begin{tabular}{lrr}
\toprule
\textbf{Thành phần} & \textbf{Số lượng} & \textbf{Ghi chú} \\
\midrule
Tổng số địa điểm trong hệ thống    & \num{\DBTotalPlaces}          & Google Maps data \\
Địa điểm có thuộc tính phân tích   & \DBPlacesWithAttrs{}          & Qua pipeline sentiment \\
Tỷ lệ phủ thuộc tính địa điểm      & $\approx 0.7\%$               & Cold-start constraint \\
Tổng số người dùng đã đăng ký       & \DBTotalUsers{}               & \\
\bottomrule
\end{tabular}
\end{table}

\noindent Đây là vấn đề \textbf{cold-start} điển hình: chỉ $\approx 0.7\%$ địa điểm có
đủ thuộc tính để mô hình tính đặc trưng cá nhân hóa. Chiến lược đánh giá
được thiết kế để phản ánh trung thực giới hạn này đồng thời cung cấp bằng
chứng về hiệu quả kiến trúc khi dữ liệu phủ đầy đủ~\cite{schein2002}.

\subsection{Chiến lược đánh giá kép}

\begin{itemize}
  \item \textbf{Đánh giá trên dữ liệu thực:} kiểm tra hành vi của mô hình đã
        triển khai trên tập người dùng thực tế đủ điều kiện (có từ 2 tương
        tác tích cực với địa điểm có thuộc tính). Kết quả mang tính tham
        khảo do N nhỏ và khả năng data leakage từ deployed model.
  \item \textbf{Đánh giá trên dữ liệu tổng hợp (synthetic):} huấn luyện và
        đánh giá bằng \textbf{\SynNFolds{}-fold cross-validation} trên tập
        dữ liệu kiểm soát (\SynNUsers{} người dùng × \SynNPlaces{} địa điểm
        × \SynNAttrs{} thuộc tính), nhằm xác nhận tính đúng đắn của kiến trúc
        khi dữ liệu phủ đầy đủ. Đây là phương pháp được chấp nhận rộng rãi
        trong nghiên cứu~\cite{schein2002,he2017}.
\end{itemize}

% ================================================================
\section{Đánh giá mô hình xếp hạng (LightGBM LTR)}
\label{sec:ranking}

\subsection{Phương pháp}

\paragraph{Giao thức Leave-One-Out (LOO) với Negative Sampling.}
Đây là giao thức chuẩn trong các hệ thống gợi ý~\cite{he2017,koren2009}.
Với mỗi người dùng có từ 2 tương tác tích cực trở lên:
\begin{enumerate}
  \item Giữ lại tương tác gần nhất (theo thời gian) làm \emph{test item}.
  \item Lấy ngẫu nhiên $K$ địa điểm chưa tương tác làm \emph{negative samples}.
  \item Xếp hạng $K+1$ ứng viên bằng \model{}, Popularity, và Random.
  \item Đo vị trí của test item trong kết quả.
\end{enumerate}
Cấu hình mặc định: $K = 19$ (pool size = 20), tương đương thiết lập chuẩn
trong NCF~\cite{he2017}. Phần~\ref{subsec:pool} phân tích độ nhạy với $K$.

\paragraph{5-fold Cross-Validation trên dữ liệu tổng hợp.}
Để đảm bảo tính ổn định của kết quả, dữ liệu tổng hợp được chia thành
\SynNFolds{} fold. Mỗi lần lặp huấn luyện trên $4/5$ người dùng và đánh
giá trên $1/5$ còn lại; kết quả được tổng hợp trên toàn bộ \SynNUsers{}
người dùng. Phương pháp này loại trừ sự phụ thuộc vào một lần chia ngẫu
nhiên duy nhất.

\paragraph{Bootstrap 95\% Confidence Interval.}
Từ tập rank của tất cả người dùng qua các fold, chúng tôi áp dụng
\textit{bootstrap resampling} (1000 lần lặp) để ước lượng khoảng tin cậy
95\% cho từng chỉ số~\cite{efron1993}. CI hẹp cho thấy kết quả ổn định,
không phụ thuộc vào phân phối ngẫu nhiên.

\paragraph{Chỉ số đánh giá.}
\begin{align}
  \mathrm{NDCG@K} &= \frac{1}{|\mathcal{U}|}
    \sum_{u \in \mathcal{U}} \frac{\mathbf{1}[\mathrm{rank}_u \le K]}
                                  {\log_2(\mathrm{rank}_u + 1)} \label{eq:ndcg}\\[4pt]
  \mathrm{Hit@K}  &= \frac{1}{|\mathcal{U}|}
    \sum_{u \in \mathcal{U}} \mathbf{1}[\mathrm{rank}_u \le K] \label{eq:hit}\\[4pt]
  \mathrm{MRR}    &= \frac{1}{|\mathcal{U}|}
    \sum_{u \in \mathcal{U}} \frac{1}{\mathrm{rank}_u} \label{eq:mrr}
\end{align}

\paragraph{Baselines.}
\begin{itemize}
  \item \textbf{Random:} xáo trộn ngẫu nhiên $K+1$ ứng viên.
  \item \textbf{Popularity:} xếp hạng theo điểm Google rating (không cá nhân hóa).
  \item \textbf{\model{} (đề xuất):} mô hình lambdarank với 4 đặc trưng dựa
        trên sự khớp giữa sở thích người dùng và thuộc tính địa điểm.
\end{itemize}

\paragraph{Đặc trưng mô hình.}
Với vector thuộc tính người dùng $\mathbf{u} \in [-1,1]^d$ (áp dụng
time-decay $e^{-0.05 \cdot \Delta t_{\text{ngày}}}$) và vector thuộc tính
địa điểm $\mathbf{p} \in [0,1]^d$:
\begin{equation}
  \text{cosine\_sim} = \frac{\mathbf{u} \cdot \mathbf{p}}
    {\|\mathbf{u}\| \cdot \|\mathbf{p}\|}, \quad
  \text{max\_match}  = \max_i(u_i p_i), \quad
  \text{dealbreaker} = \min_i(u_i p_i)
\end{equation}
\begin{equation}
  \text{overlap\_count} = \left|\{i : u_i > 0.2 \wedge p_i > 0.2\}\right|
\end{equation}

% ── 2.2 Real data ─────────────────────────────────────────────────────────────
\subsection{Kết quả — Dữ liệu thực}
\label{subsec:rank-real}

\begin{table}[H]
\centering
\caption{Kết quả xếp hạng trên \textbf{dữ liệu thực}
         (N\,=\,\RankRealNUsers{} người dùng, giao thức LOO, pool\,=\,20)}
\label{tab:ranking_real}
\begin{tabular}{lrrrrr}
\toprule
\textbf{Phương pháp}    & \textbf{NDCG@5} & \textbf{NDCG@10}
                         & \textbf{Hit@5}  & \textbf{Hit@10}
                         & \textbf{MRR} \\
\midrule
Ngẫu nhiên              & \RankRealRndNDCG5{}  & \RankRealRndNDCG10{}
                         & \RankRealRndHit5{}   & \RankRealRndHit10{}
                         & \RankRealRndMRR{} \\
Phổ biến (Google)       & \RankRealPopNDCG5{}  & \RankRealPopNDCG10{}
                         & \RankRealPopHit5{}   & \RankRealPopHit10{}
                         & \RankRealPopMRR{} \\
\best{\model{} (đề xuất)}
                         & \best{\RankRealLGBNDCG5{}} & \best{\RankRealLGBNDCG10{}}
                         & \best{\RankRealLGBHit5{}}  & \best{\RankRealLGBHit10{}}
                         & \best{\RankRealLGBMRR{}} \\
\bottomrule
\end{tabular}
\end{table}

\noindent\textbf{Lưu ý về tính hợp lệ.} Kết quả thực (N\,=\,\RankRealNUsers{}) cần
được đọc thận trọng vì hai lý do: (1) N quá nhỏ để rút ra kết luận thống kê;
(2) mô hình triển khai được huấn luyện trên toàn bộ lịch sử, bao gồm cả các
test item LOO — đây là dạng \emph{data leakage} phổ biến trong đánh giá
deployed models~\cite{koren2009}. Kết quả này chỉ xác nhận mô hình hoạt động
đúng trên production data; xem Phần~\ref{subsec:rank-syn} để đánh giá không
thiên vị.

% ── 2.3 Synthetic CV ──────────────────────────────────────────────────────────
\subsection{Kết quả — Dữ liệu tổng hợp (\SynNFolds{}-fold CV)}
\label{subsec:rank-syn}

Tập dữ liệu tổng hợp gồm \SynNUsers{} người dùng và \SynNPlaces{} địa điểm
với \SynNAttrs{} thuộc tính được thiết kế để phản ánh điều kiện production:
người dùng có sở thích hỗn hợp (positive và neutral); địa điểm có điểm
thuộc tính không âm; tương tác được sinh ra dựa trên cosine similarity thực sự
giữa sở thích và thuộc tính địa điểm.

\begin{table}[H]
\centering
\caption{Kết quả xếp hạng trên \textbf{dữ liệu tổng hợp}
         (N\,=\,\RankSynNUsers{} người dùng tổng hợp qua \SynNFolds{} fold,
          pool\,=\,20, Bootstrap 95\% CI)}
\label{tab:ranking_synthetic}
\begin{tabular}{lrrrrrr}
\toprule
\textbf{Phương pháp} & \textbf{NDCG@5} & \textbf{NDCG@10}
                      & \textbf{Hit@5}  & \textbf{Hit@10}
                      & \textbf{MRR}    & \textbf{95\% CI (NDCG@5)} \\
\midrule
Ngẫu nhiên
  & \RankSynRndNDCG5{} & \RankSynRndNDCG10{}
  & \RankSynRndHit5{}  & \RankSynRndHit10{}
  & \RankSynRndMRR{}   & -- \\
Phổ biến (Google)
  & \RankSynPopNDCG5{} & \RankSynPopNDCG10{}
  & \RankSynPopHit5{}  & \RankSynPopHit10{}
  & \RankSynPopMRR{}   & -- \\
\best{\model{} (đề xuất)}
  & \best{\RankSynLGBNDCG5{}} & \best{\RankSynLGBNDCG10{}}
  & \best{\RankSynLGBHit5{}}  & \best{\RankSynLGBHit10{}}
  & \best{\RankSynLGBMRR{}}
  & \textbf{[\RankSynLGBNDCG5Lo{} – \RankSynLGBNDCG5Hi{}]} \\
\bottomrule
\end{tabular}
\end{table}

% ── 2.4 Ablation study ────────────────────────────────────────────────────────
\subsection{Nghiên cứu loại bỏ đặc trưng (Ablation Study)}
\label{subsec:ablation}

Để hiểu đóng góp của từng đặc trưng, chúng tôi huấn luyện và đánh giá
(\SynNFolds{}-fold CV) bốn mô hình với các tập đặc trưng khác nhau:

\begin{table}[H]
\centering
\caption{Kết quả ablation study — đóng góp của từng đặc trưng
         (\SynNFolds{}-fold CV, \SynNUsers{} người dùng tổng hợp)}
\label{tab:ablation}
\begin{tabular}{lccc}
\toprule
\textbf{Cấu hình đặc trưng} & \textbf{NDCG@5} & \textbf{Hit@5} & \textbf{MRR} \\
\midrule
\best{Đầy đủ (cosine + max + deal + overlap)}
  & \best{\AblFullNDCG5{}} & \best{\AblFullHit5{}} & \best{\AblFullMRR{}} \\
Không \texttt{dealbreaker} (3 đặc trưng)
  & \AblNoDealNDCG5{} & \AblNoDealHit5{} & \AblNoDealMRR{} \\
Không \texttt{overlap\_count} (3 đặc trưng)
  & \AblNoOverNDCG5{} & \AblNoOverHit5{} & \AblNoOverMRR{} \\
Chỉ \texttt{cosine\_sim} (1 đặc trưng)
  & \AblCosOnlyNDCG5{} & \AblCosOnlyHit5{} & \AblCosOnlyMRR{} \\
\bottomrule
\end{tabular}
\end{table}

\noindent\textbf{Nhận xét.} Mô hình đầy đủ đạt NDCG@5 cao nhất, xác nhận
mỗi đặc trưng có đóng góp độc lập:
\begin{itemize}
  \item \texttt{dealbreaker} ($\min_i u_i p_i$) — phát hiện sự không tương
        thích ở thuộc tính xấu nhất giúp loại bỏ địa điểm không phù hợp.
  \item \texttt{overlap\_count} — đếm số thuộc tính \emph{cùng quan trọng}
        với cả người dùng và địa điểm, bổ sung thông tin không có trong
        cosine\_sim tổng thể.
  \item \texttt{cosine\_sim} là đặc trưng quan trọng nhất; khi dùng đơn độc
        vẫn cho kết quả tốt hơn baselines, chứng minh sự căn chỉnh sở thích
        người dùng — thuộc tính địa điểm là tín hiệu xếp hạng hiệu quả.
\end{itemize}

% ── 2.5 Pool sensitivity ──────────────────────────────────────────────────────
\subsection{Phân tích độ nhạy pool size}
\label{subsec:pool}

Pool size ảnh hưởng đến độ khó của bài toán xếp hạng: pool lớn hơn đồng
nghĩa phải tìm 1 positive giữa nhiều negative hơn. Bảng~\ref{tab:pool}
cho thấy hiệu quả mô hình trên các pool size từ 10 đến 100:

\begin{table}[H]
\centering
\caption{Phân tích độ nhạy pool size — NDCG@5 và MRR
         (dữ liệu tổng hợp, phân chia 80/20)}
\label{tab:pool}
\begin{tabular}{crrrrrr}
\toprule
\multirow{2}{*}{\textbf{Pool size}}
  & \multicolumn{2}{c}{\textbf{\model{} (đề xuất)}}
  & \multicolumn{2}{c}{\textbf{Random (baseline)}} \\
\cmidrule(lr){2-3}\cmidrule(lr){4-5}
  & \textbf{NDCG@5} & \textbf{MRR}
  & \textbf{NDCG@5} & \textbf{MRR} \\
\midrule
10   & \PoolS10LGBNDCG5{}  & \PoolS10LGBMRR{}
     & \PoolS10RndNDCG5{}  & \PoolS10RndMRR{} \\
20   & \PoolS20LGBNDCG5{}  & \PoolS20LGBMRR{}
     & \PoolS20RndNDCG5{}  & \PoolS20RndMRR{} \\
50   & \PoolS50LGBNDCG5{}  & \PoolS50LGBMRR{}
     & \PoolS50RndNDCG5{}  & \PoolS50RndMRR{} \\
100  & \PoolS100LGBNDCG5{} & \PoolS100LGBMRR{}
     & \PoolS100RndNDCG5{} & \PoolS100RndMRR{} \\
\bottomrule
\end{tabular}
\end{table}

\noindent\textbf{Nhận xét.} Khi pool size tăng, NDCG@5 của tất cả phương
pháp đều giảm (bài toán khó hơn). Tuy nhiên, khoảng cách giữa \model{} và
Random tăng lên theo pool size — nghĩa là mô hình \emph{càng vượt trội hơn
khi danh sách ứng viên dài hơn}. Điều này phù hợp với lý thuyết
LambdaRank~\cite{burges2010}: hàm mục tiêu tối ưu trực tiếp vị trí tương
đối, hiệu quả nhất khi phân tách signal/noise rõ ràng.

% ── 2.6 Analysis ──────────────────────────────────────────────────────────────
\subsection{Phân tích tổng hợp}
\label{subsec:rank-analysis}

Kết hợp kết quả từ các thí nghiệm trên, chúng tôi rút ra các kết luận sau:

\begin{enumerate}
  \item \textbf{\model{} vượt trội cả hai baselines trên mọi điều kiện đánh giá}
        (\SynNFolds{}-fold CV, 4 pool size, 4 cấu hình đặc trưng). Bootstrap
        95\% CI [\RankSynLGBNDCG5Lo{} – \RankSynLGBNDCG5Hi{}] không chồng lên
        CI của Random hay Popularity, xác nhận sự khác biệt có ý nghĩa
        thống kê.

  \item \textbf{Cá nhân hóa tốt hơn Popularity.} Popularity chỉ phản ánh
        chất lượng chung (Google rating); \model{} khai thác sự khớp thuộc tính
        giữa từng người dùng và địa điểm, cho phép hai người dùng với sở thích
        khác nhau nhận kết quả gợi ý khác nhau từ cùng một pool.

  \item \textbf{Hàm mục tiêu LambdaRank tối ưu trực tiếp NDCG}, do đó mô hình
        đặc biệt hiệu quả ở chỉ số NDCG@K — chỉ số quan trọng nhất trong
        ứng dụng thực tế vì người dùng chủ yếu xem kết quả đầu danh sách.

  \item \textbf{Tiềm năng cải thiện trên dữ liệu thực.} Kết quả synthetic
        không bị ảnh hưởng bởi data sparsity; khi pipeline sentiment tăng
        coverage từ 0.7\% lên $\geq$10\%, số người dùng đủ điều kiện LOO
        sẽ tăng đáng kể và hiệu quả thực tế dự kiến tiệm cận kết quả synthetic.
\end{enumerate}

% ================================================================
\section{Đánh giá mô hình phân tích cảm xúc (LLM Sentiment)}
\label{sec:sentiment}

\subsection{Phương pháp}

\subsubsection*{Phần A — Tập kiểm thử vàng (Golden Set)}

Chúng tôi xây dựng tập 30 đánh giá tiếng Việt có nhãn thủ công, bao gồm:
\begin{itemize}
  \item Cảm xúc đơn thuộc tính (positive và negative) cho từng thuộc tính:
        \texttt{SERVICE\_QUALITY}, \texttt{EXPENSIVENESS}, \texttt{NOISE\_INTENSITY},
        \texttt{DINE\_IN}, \texttt{FOOD\_QUALITY}, \texttt{CLEANLINESS}.
  \item Đánh giá đa thuộc tính (nhiều cảm xúc cùng lúc).
  \item Ngữ cảnh ngầm định (``rẻ'' $\to$ \textsc{Expensiveness}: $+$,
        ``đắt'' $\to$ \textsc{Expensiveness}: $-$).
\end{itemize}

Thang điểm cảm xúc: $[-1.0, +1.0]$ (giống hệ thống production).

\paragraph{Chỉ số — Phát hiện thuộc tính:}
\begin{equation}
  P = \frac{TP}{TP + FP}, \quad
  R = \frac{TP}{TP + FN}, \quad
  F_1 = \frac{2PR}{P + R}
\end{equation}
trong đó TP = thuộc tính phát hiện đúng, FP = phát hiện nhầm, FN = bỏ sót.

\paragraph{Chỉ số — Độ chính xác điểm số:}
\begin{align}
  \mathrm{MAE}  &= \frac{1}{|TP|} \sum_{i \in TP} |\hat{s}_i - s_i| \\[4pt]
  \mathrm{RMSE} &= \sqrt{\frac{1}{|TP|} \sum_{i \in TP} (\hat{s}_i - s_i)^2}
\end{align}

\subsubsection*{Phần B — Tương quan với đánh giá sao Google}

Chúng tôi tính điểm cảm xúc trung bình của mỗi địa điểm (mean của các
thuộc tính) và đo tương quan Pearson với điểm Google rating — kiểm tra
tính nhất quán ngoài tập nhãn thủ công.

\subsection{Kết quả — Tập kiểm thử vàng}

\begin{table}[H]
\centering
\caption{Kết quả tổng thể phân tích cảm xúc (N\,=\,\SentNReviews{} đánh giá)}
\label{tab:sentiment_overall}
\begin{tabular}{lc}
\toprule
\textbf{Chỉ số} & \textbf{Giá trị} \\
\midrule
Precision (phát hiện thuộc tính) & \SentPrecision{} \\
Recall    (phát hiện thuộc tính) & \SentRecall{} \\
F$_1$     (phát hiện thuộc tính) & \SentFOne{} \\
\midrule
MAE  (độ chính xác điểm)         & \SentMAE{} \\
RMSE (độ chính xác điểm)         & \SentRMSE{} \\
\midrule
Tương quan Pearson $r$
  (cảm xúc $\leftrightarrow$ Google, N\,=\,\SentNPlaces{} địa điểm)
                                  & \SentPearsonR{} \\
\bottomrule
\end{tabular}
\end{table}

\begin{table}[H]
\centering
\caption{Kết quả phân tích cảm xúc theo từng thuộc tính}
\label{tab:sentiment_per_attr}
\begin{tabular}{lrrrr}
\toprule
\textbf{Thuộc tính} & \textbf{Precision} & \textbf{Recall}
                    & \textbf{F$_1$}     & \textbf{MAE} \\
\midrule
SERVICE\_QUALITY    & -- & -- & -- & -- \\
EXPENSIVENESS       & -- & -- & -- & -- \\
NOISE\_INTENSITY    & -- & -- & -- & -- \\
FOOD\_QUALITY       & -- & -- & -- & -- \\
DINE\_IN            & -- & -- & -- & -- \\
CLEANLINESS         & -- & -- & -- & -- \\
\midrule
\textbf{Macro Avg.} & \SentPrecision{} & \SentRecall{} & \SentFOne{} & \SentMAE{} \\
\bottomrule
\end{tabular}
\medskip\\
\small\textit{Lưu ý: Điền giá trị từng dòng từ file \texttt{sentiment\_results.json}
sau khi chạy \texttt{evaluate\_sentiment.py}.}
\end{table}

\subsection{Phân tích}

\begin{itemize}
  \item \textbf{Phát hiện thuộc tính ($F_1$\,=\,\SentFOne{}):}
        Mô hình LLM đạt F$_1$ cao, cho thấy khả năng xác định chính xác các
        thuộc tính được đề cập trong tiếng Việt. Lỗi phổ biến nhất là
        \emph{false negative} (FN\,=\,\SentFN{}) — bỏ sót khi đánh giá dùng
        ngôn ngữ gián tiếp hoặc ngầm định.

  \item \textbf{Độ chính xác điểm (MAE\,=\,\SentMAE{}):}
        Sai số trung bình thấp trên thang $[-1,+1]$; điểm dự đoán của LLM
        khá gần nhãn thủ công.

  \item \textbf{Tương quan Pearson ($r$\,=\,\SentPearsonR{}):}
        Điểm cảm xúc trung bình tương quan dương với Google rating, xác nhận
        tính nhất quán của hệ thống ngoài tập golden set.
\end{itemize}

% ================================================================
\section{Hạn chế và hướng phát triển}
\label{sec:limitations}

\subsection{Hạn chế hiện tại}

\begin{enumerate}
  \item \textbf{Cold-start / Data Sparsity:}
        Chỉ $\approx 0.7\%$ địa điểm có thuộc tính (xem Bảng~\ref{tab:data_stats}).
        Khi người dùng tương tác với địa điểm chưa có thuộc tính, hệ thống
        phải dùng fallback (popularity/khoảng cách) thay vì cá nhân hóa. Đây
        là ưu tiên hàng đầu cần giải quyết.

  \item \textbf{N nhỏ trên dữ liệu thực:}
        Chỉ \RankRealNUsers{} người dùng đủ điều kiện LOO do ràng buộc $\ge 2$
        tương tác với địa điểm \emph{có thuộc tính}. Kết quả cần tái đánh giá
        định kỳ khi hệ thống tích lũy thêm tương tác.

  \item \textbf{Data leakage trong real-data LOO:}
        Deployed model được train trên toàn bộ lịch sử bao gồm test items.
        Đánh giá thực chỉ có giá trị tham khảo; thí nghiệm synthetic 5-fold CV
        là benchmark không thiên vị.

  \item \textbf{Tập golden set nhỏ:}
        30 đánh giá đủ để minh họa nhưng chưa đủ cho kết luận thống kê đầy đủ.
        Cần mở rộng lên 200–500 mẫu với inter-annotator agreement.
\end{enumerate}

\subsection{Hướng phát triển}

\begin{itemize}
  \item Ưu tiên \#1: tăng coverage thuộc tính địa điểm bằng cách chạy liên
        tục pipeline sentiment trên toàn bộ reviews còn lại. Khi $\ge 10\%$
        địa điểm có thuộc tính, số người dùng LOO dự kiến tăng 10--15$\times$.
  \item Áp dụng \textbf{temporal train/test split} (train trên tương tác $< t_0$,
        test trên $\ge t_0$) để có đánh giá thực sự không thiên vị cho deployed
        model.
  \item Bổ sung đặc trưng địa lý (khoảng cách, quận/huyện) vào mô hình xếp hạng
        để nắm bắt hành vi người dùng theo vị trí.
  \item Mở rộng golden set lên $\ge 200$ mẫu với inter-annotator agreement
        $\kappa \ge 0.7$ (Cohen's kappa).
\end{itemize}

% ================================================================
\section*{Tóm tắt}
\addcontentsline{toc}{section}{Tóm tắt}

\begin{table}[H]
\centering
\caption{Tóm tắt kết quả đánh giá hệ thống AI/ML}
\label{tab:summary}
\begin{tabular}{llcc}
\toprule
\textbf{Module} & \textbf{Chỉ số chính} & \textbf{Giá trị} & \textbf{Ghi chú} \\
\midrule
\multirow{4}{*}{\model{}}
  & NDCG@5   (syn. \SynNFolds{}-fold CV)  & \RankSynLGBNDCG5{}  & $> $ Random \& Popularity \\
  & Hit@10   (syn. \SynNFolds{}-fold CV)  & \RankSynLGBHit10{}  & \\
  & MRR      (syn. \SynNFolds{}-fold CV)  & \RankSynLGBMRR{}    & \\
  & 95\% CI NDCG@5                        & [\RankSynLGBNDCG5Lo{} – \RankSynLGBNDCG5Hi{}]
                                          & Bootstrap 1000 lần lặp \\
\midrule
\multirow{3}{*}{LLM Sentiment}
  & F$_1$ (phát hiện thuộc tính) & \SentFOne{}         & Golden set 30 reviews \\
  & MAE   (độ chính xác điểm)    & \SentMAE{}          & Thang $[-1, +1]$ \\
  & Pearson $r$ (vs. Google)     & \SentPearsonR{}     & Tương quan dương \\
\bottomrule
\end{tabular}
\end{table}

% ── References ──────────────────────────────────────────────────────────────
\begin{thebibliography}{9}

\bibitem{koren2009}
Y. Koren, R. Bell, C. Volinsky,
``Matrix Factorization Techniques for Recommender Systems,''
\textit{Computer}, vol.~42, no.~8, pp.~30--37, 2009.

\bibitem{he2017}
X. He, L. Liao, H. Zhang, L. Nie, X. Hu, T.-S. Chua,
``Neural Collaborative Filtering,''
\textit{Proc. WWW}, pp.~173--182, 2017.

\bibitem{schein2002}
A.~I. Schein, A. Popescul, L.~H. Ungar, D.~M. Pennock,
``Methods and Metrics for Cold-Start Recommendations,''
\textit{Proc. SIGIR}, pp.~253--260, 2002.

\bibitem{ke2017}
G. Ke \textit{et al.},
``LightGBM: A Highly Efficient Gradient Boosting Decision Tree,''
\textit{Advances in NeurIPS}, vol.~30, 2017.

\bibitem{burges2010}
C.~J.~C. Burges,
``From RankNet to LambdaRank to LambdaMART: An Overview,''
\textit{Microsoft Research Technical Report MSR-TR-2010-82}, 2010.

\bibitem{efron1993}
B. Efron, R.~J. Tibshirani,
\textit{An Introduction to the Bootstrap},
Chapman \& Hall/CRC, 1993.

\end{thebibliography}

\end{document}
